% How accurate LIGHT v0_2 is: its accuracy, its system map and its reaction
% rates (Appendix A). The v0 and v0_1 overviews describe those versions. No
% number here is typed by hand. They come from five generated files:
%   system_map_v0_2_numbers.tex    make_system_map_v0_2.py
%   overview_v0_2_rates.tex, overview_v0_2_fast.tex, overview_v0_2_slow.tex
%                                  make_overview_v0_2.py
%   energy_ladder_v0_2.tex         make_energy_ladder.py
% Build from LIGHT/docs:
%   python3 make_system_map_v0_2.py && python3 make_overview_v0_2.py
%   python3 make_energy_ladder.py
%   latexmk -pdf system_map_v0_2.tex && latexmk -pdf accuracy_overview_v0_2.tex
\documentclass[11pt]{article}
\usepackage[a4paper,margin=2.2cm]{geometry}
\usepackage[T1]{fontenc}
\usepackage[scaled=0.95]{helvet}
\renewcommand{\familydefault}{\sfdefault}
\usepackage{amsmath}
\usepackage{sansmath}
\usepackage{booktabs,tabularx,array,longtable}
\usepackage{xcolor}
\usepackage{pdfpages}
\usepackage{float}
\usepackage{enumitem}
\usepackage{needspace}
\usepackage{tikz}
\usetikzlibrary{arrows.meta,decorations.pathmorphing}
\setlist{nosep, leftmargin=1.2em}
\setlength{\parindent}{0pt}
\setlength{\parskip}{5pt}
\definecolor{ink2}{HTML}{52514E}

% LIGHT v0_2 map numbers. GENERATED by make_system_map_v0_2.py
% from the model (0.1.2). Do not edit. Rerun the script instead.
\providecommand\wmdef[2]{\expandafter\def\csname wm@#1\endcsname{#2}}
\wmdef{e.P}{100.0}
\wmdef{e.pen}{11.0}
\wmdef{e.pen2}{0.0}
\wmdef{e.G1}{85.9}
\wmdef{e.C1}{12.4}
\wmdef{e.C2}{0.4}
\wmdef{e.F2}{2.8}
\wmdef{e.G6}{6.2}
\wmdef{e.dU}{0.0}
\wmdef{e.cap}{3.3}
\wmdef{C1.kr}{10.6}
\wmdef{C1.heat}{1.8}
\wmdef{F2e.F}{1.4}
\wmdef{F2e.heat}{1.9}
\wmdef{F2e.anion}{-0.6}
\wmdef{krs.in}{113.8}
\wmdef{krs.harp}{47.2}
\wmdef{krs.G2}{37.4}
\wmdef{krs.AK8}{29.2}
\wmdef{krs.left}{0.0}
\wmdef{brs.harp}{41.5}
\wmdef{brs.G2}{32.8}
\wmdef{brs.AK8}{25.6}
\wmdef{harp.B}{15.6}
\wmdef{harp.C}{9.5}
\wmdef{harp.F}{4.2}
\wmdef{harp.heat}{18.0}
\wmdef{G2.ion}{26.4}
\wmdef{C2.ion}{0.4}
\wmdef{AK8.heat}{2.1}
\wmdef{k2s.in}{27.1}
\wmdef{k2s.RAD}{21.8}
\wmdef{k2s.F2}{4.7}
\wmdef{k2s.F}{0.0}
\wmdef{k2s.pen}{0.0}
\wmdef{k2s.left}{0.5}
\wmdef{brk2.RAD}{80.8}
\wmdef{brk2.F2}{17.4}
\wmdef{RAD.hv}{20.0}
\wmdef{RAD.heat}{1.8}
\wmdef{NK7.B}{1.5}
\wmdef{BR.k2f}{1.3}
\wmdef{K2F2.F}{0.4}
\wmdef{K2F2.heat}{1.4}
\wmdef{k2F.B}{0.0}
\wmdef{k2F.Fin}{0.0}
\wmdef{k2F.heat}{0.0}
\wmdef{pen2.ion}{0.0}
\wmdef{K15.heat}{2.2}
\wmdef{R11p.in}{22.4}
\wmdef{R11p.kr}{17.3}
\wmdef{R11p.heat}{5.1}
\wmdef{ii.in}{3.0}
\wmdef{ii.anion}{-0.6}
\wmdef{ii.B}{1.3}
\wmdef{ii.heat}{1.1}
\wmdef{ions.left}{-0.7}
\wmdef{anion.left}{0.0}
\wmdef{mix.net}{2.7}
\wmdef{mix.netC}{2.7}
\wmdef{mix.heat}{0.0}
\wmdef{mix.BC}{10.1}
\wmdef{mix.CB}{12.8}
\wmdef{B.G4}{5.5}
\wmdef{B.G4A}{0.2}
\wmdef{B.G5}{13.8}
\wmdef{B.Q}{1.6}
\wmdef{B.QF2}{0.9}
\wmdef{B.QKR}{0.0}
\wmdef{B.QE}{0.7}
\wmdef{B.left}{0.0}
\wmdef{C.G4C}{0.9}
\wmdef{C.G5C}{3.5}
\wmdef{C.Q}{2.4}
\wmdef{C.QF2}{0.7}
\wmdef{C.QKR}{0.9}
\wmdef{C.QE}{0.8}
\wmdef{C.left}{0.0}
\wmdef{G4.hv}{4.7}
\wmdef{G4.F}{0.7}
\wmdef{G4.heat}{0.1}
\wmdef{G4A.hv}{0.1}
\wmdef{G4A.F}{0.0}
\wmdef{G4A.heat}{0.0}
\wmdef{G4C.hv}{0.7}
\wmdef{G4C.F}{0.1}
\wmdef{G4C.heat}{0.1}
\wmdef{QB.F}{0.5}
\wmdef{QB.heat}{1.2}
\wmdef{QC.F}{0.5}
\wmdef{QC.heat}{1.9}
\wmdef{tri.k2f}{15.6}
\wmdef{tri.heat}{1.7}
\wmdef{k2f.C3}{8.7}
\wmdef{k2f.Q2F}{8.2}
\wmdef{k2f.Q2FF}{0.0}
\wmdef{k2f.left}{0.0}
\wmdef{brk.C3}{51.4}
\wmdef{brk.Q2F}{48.5}
\wmdef{C3.hv}{4.9}
\wmdef{C3.F}{1.3}
\wmdef{C3.heat}{2.5}
\wmdef{Q2F.F}{1.2}
\wmdef{Q2F.heat}{7.0}
\wmdef{Q2FF.F}{0.0}
\wmdef{Q2FF.heat}{0.0}
\wmdef{FF.in}{0.1}
\wmdef{FF.heat}{0.1}
\wmdef{out.248}{4.7}
\wmdef{out.272}{0.1}
\wmdef{out.275}{0.7}
\wmdef{out.420}{4.9}
\wmdef{out.147}{20.0}
\wmdef{out.heat}{53.2}
\wmdef{out.elastic}{6.2}
\wmdef{out.ions}{-0.7}
\wmdef{out.F}{10.4}
\wmdef{out.other}{0.5}
\wmdef{G4x.F}{0.8}
\wmdef{G4x.heat}{0.1}
\wmdef{iiNK7.B}{2.8}
\wmdef{k2s.fates}{26.6}
\wmdef{KRx.heat}{3.3}
\wmdef{version}{0.1.2}
\wmdef{p}{1}
\wmdef{xF2}{0.5}
\wmdef{Tg0}{300}
\wmdef{P}{300}
\wmdef{pulse}{1}
\wmdef{tend}{2}
\wmdef{ne}{10^{15}}
\wmdef{neF2}{1.68{\times}10^{14}}
\wmdef{nePeak}{1.70{\times}10^{15}}
\wmdef{N}{2.45{\times}10^{19}}
\wmdef{Edep}{0.30}
\wmdef{Te}{4.45}
\wmdef{TeMid}{3.81}
\wmdef{rTail}{0.16}
\wmdef{TtMid}{0.59}
\wmdef{TgEnd}{649}
\wmdef{TgMax}{649}
\wmdef{pEnd}{2.18}
\wmdef{F2burn}{640}
\wmdef{F2left}{0.20}
\wmdef{krsFormed}{2.15{\times}10^{17}}
\wmdef{k2sFormed}{5.51{\times}10^{16}}
\wmdef{krfFormed}{8.93{\times}10^{16}}
\wmdef{k2fFormed}{6.01{\times}10^{16}}
\wmdef{krf.eff}{23.8}
\wmdef{frm.ii}{4.7}
\wmdef{frm.kr2}{5.4}
\wmdef{brf.G4}{19.6}
\wmdef{brf.G4A}{0.7}
\wmdef{brf.G4C}{3.3}
\wmdef{brf.tri}{62.0}
\wmdef{brf.e}{5.4}
\wmdef{brf.F2}{5.7}
\wmdef{brf.Kr}{3.4}
\wmdef{brf.QB}{5.8}
\wmdef{brf.QC}{8.7}
\wmdef{ion.R11p}{88.7}
\wmdef{ion.ii}{11.3}
\wmdef{brs.AK8.F2}{15.1}
\wmdef{k.harp}{8.47{\times}10^{-10}}
\wmdef{pc.harpY}{90}
\wmdef{pc.harpB}{56}
\wmdef{pc.harpC}{34}
\wmdef{pc.harpX}{10}
\wmdef{pc.fC}{38}
\wmdef{k.G2}{5.0{\times}10^{-10}}
\wmdef{tauB}{6.99}
\wmdef{pc.BA}{3.5}
\wmdef{tauC}{50}
\wmdef{tauC3}{220}
\wmdef{gapBC}{120}
\wmdef{k.BC}{1.05{\times}10^{-11}}
\wmdef{Keq}{1.10}
\wmdef{k.G5}{1.9{\times}10^{-30}}
\wmdef{k.G5C}{3.79{\times}10^{-31}}
\wmdef{tpow.G5}{-3.32}
\wmdef{k.Q2F}{1.2{\times}10^{-10}}
\wmdef{k.K15}{2.4{\times}10^{-31}}
\wmdef{tpow.K15}{-0.75}
\wmdef{k.R11p}{1.6{\times}10^{-6}}
\wmdef{pow.R11p}{0.55}
\wmdef{cap.R11p}{0.95}
\wmdef{k.FF}{2.8{\times}10^{-34}}
\wmdef{n3.ii}{-3}
\wmdef{k.AK8}{1.7{\times}10^{-32}}
\wmdef{tpow.AK8}{-0.64}
\wmdef{tauRAD}{270}
\wmdef{k.kr2F2}{2.1{\times}10^{-10}}
\wmdef{pc.NK7}{50}
\wmdef{E.kr2s}{9.20}
\wmdef{E.kr2s.bind}{0.715}
\wmdef{hv.RAD}{8.43}
\wmdef{E.anion}{2.6}
\wmdef{hv.G4C}{4.51}
\wmdef{hv.G4}{5.00}
\wmdef{hv.G4A}{4.56}
\wmdef{hv.C3}{2.95}
\wmdef{E.B}{5.856}
\wmdef{notes.krs}{98.3}
\wmdef{notes.krp}{26.8}
\wmdef{notes.emit}{30.4}
\wmdef{nodes}{43}
\wmdef{om.electron}{4.4}
\wmdef{om.basic}{4.7}
\wmdef{om.v0}{16.5}
\wmdef{bd.kr2_f.def}{0}
\wmdef{bd.kr2_f.lo}{0}
\wmdef{bd.kr2_f.hi}{1.0{\times}10^{-9}}
\wmdef{bd.kr2_f.lo248}{4.67}
\wmdef{bd.kr2_f.hi248}{9.50}
\wmdef{bd.trimer_tg_pow.def}{-3.32}
\wmdef{bd.trimer_tg_pow.lo}{-3.32}
\wmdef{bd.trimer_tg_pow.hi}{-0.50}
\wmdef{bd.trimer_tg_pow.lo248}{4.67}
\wmdef{bd.trimer_tg_pow.hi248}{2.22}
\wmdef{bd.kr_pen.def}{5.0{\times}10^{-10}}
\wmdef{bd.kr_pen.lo}{1.5{\times}10^{-10}}
\wmdef{bd.kr_pen.hi}{5.0{\times}10^{-10}}
\wmdef{bd.kr_pen.lo248}{4.87}
\wmdef{bd.kr_pen.hi248}{4.67}
\wmdef{bd.kr2f_f2.def}{1.2{\times}10^{-10}}
\wmdef{bd.kr2f_f2.lo}{1.2{\times}10^{-10}}
\wmdef{bd.kr2f_f2.hi}{4.3{\times}10^{-10}}
\wmdef{bd.kr2f_f2.lo248}{4.67}
\wmdef{bd.kr2f_f2.hi248}{4.88}
\wmdef{bd.kr2f_f.def}{0}
\wmdef{bd.kr2f_f.lo}{0}
\wmdef{bd.kr2f_f.hi}{1.8{\times}10^{-10}}
\wmdef{bd.kr2f_f.lo248}{4.67}
\wmdef{bd.kr2f_f.hi248}{4.87}
\wmdef{bd.kr2_form.def}{1.7{\times}10^{-32}}
\wmdef{bd.kr2_form.lo}{1.7{\times}10^{-32}}
\wmdef{bd.kr2_form.hi}{5.5{\times}10^{-32}}
\wmdef{bd.kr2_form.lo248}{4.67}
\wmdef{bd.kr2_form.hi248}{4.52}
\wmdef{bd.kr2_f2.def}{2.1{\times}10^{-10}}
\wmdef{bd.kr2_f2.lo}{2.1{\times}10^{-10}}
\wmdef{bd.kr2_f2.hi}{5.0{\times}10^{-10}}
\wmdef{bd.kr2_f2.lo248}{4.67}
\wmdef{bd.kr2_f2.hi248}{4.60}
\wmdef{bd.kr2_pen.def}{0}
\wmdef{bd.kr2_pen.lo}{0}
\wmdef{bd.kr2_pen.hi}{5.0{\times}10^{-10}}
\wmdef{bd.kr2_pen.lo248}{4.67}
\wmdef{bd.kr2_pen.hi248}{4.65}
\wmdef{bd.minor}{0.21}
\wmdef{bd.n}{8}
\wmdef{bd.nzero}{3}
\def\wmBands{\Krts{} + F \ra KrF(B) + Kr & $0$ & $0$ to $1.0{\times}10^{-9}$\unit{cm^3\,s^{-1}} & 4.67 to 9.50 & 20.0 to 0.5 & 649 to 726\\
the trimer's $(T_g/300)^n$ & $-3.32$ & $-3.32$ to $-0.50$ & 4.67 to 2.22 & 20.0 to 19.5 & 649 to 658\\
\Krs{} + \Krs{} \ra \Krp{} + Kr + \el{} & $5.0{\times}10^{-10}$ & $1.5{\times}10^{-10}$ to $5.0{\times}10^{-10}$\unit{cm^3\,s^{-1}} & 4.87 to 4.67 & 25.0 to 20.0 & 616 to 649\\
\KtF{} + \Ft{} & $1.2{\times}10^{-10}$ & $1.2{\times}10^{-10}$ to $4.3{\times}10^{-10}$\unit{cm^3\,s^{-1}} & 4.67 to 4.88 & 20.0 to 19.9 & 649 to 663\\
\KtF{} + F \ra 2Kr + 2F & $0$ & $0$ to $1.8{\times}10^{-10}$\unit{cm^3\,s^{-1}} & 4.67 to 4.87 & 20.0 to 19.9 & 649 to 670\\
\Krs{} + 2Kr \ra \Krts{} + Kr, 300\unit{K} & $1.7{\times}10^{-32}$ & $1.7{\times}10^{-32}$ to $5.5{\times}10^{-32}$\unit{cm^6\,s^{-1}} & 4.67 to 4.52 & 20.0 to 28.0 & 649 to 601\\
\Krts{} + \Ft{}, both channels & $2.1{\times}10^{-10}$ & $2.1{\times}10^{-10}$ to $5.0{\times}10^{-10}$\unit{cm^3\,s^{-1}} & 4.67 to 4.60 & 20.0 to 19.1 & 649 to 654\\
\Krts{} + \Krts{} \ra \Krtp{} + 2Kr + \el{} & $0$ & $0$ to $5.0{\times}10^{-10}$\unit{cm^3\,s^{-1}} & 4.67 to 4.65 & 20.0 to 10.8 & 649 to 704\\}

% LIGHT v0_2 reaction-rate tables. GENERATED by make_overview_v0_2.py
% from LIGHT v0_2 (0.1.2). Do not edit. Rerun the script instead.
\providecommand\widef[2]{\expandafter\def\csname wi@#1\endcsname{#2}}
\widef{rates.teMid}{3.81}
\widef{rates.rMid}{0.16}
\widef{rates.nBasic}{44}
\widef{rates.nBandRows}{8}
\widef{rates.nGroups}{2}
\widef{rates.groups}{electron}
\widef{rates.v0groups}{quench, kr2, ions, electron, bc, f2e}
\def\wirBasic{G1 & \el{} + Kr \ra \el{} + \Krs{} & $\langle\sigma v\rangle$, Table A.3 & $T_e$ & Biagi v7.1, the whole 5s manifold at 9.915 eV\\
C1a & \el{} + Kr \ra \el{} + \Krs{} & $\langle\sigma v\rangle$, Table A.3 & $T_e$ & Biagi v7.1 5p group at 11.30 eV, prompt cascade to 5s\\
C1b & \el{} + Kr \ra \el{} + \Krs{} & $\langle\sigma v\rangle$, Table A.3 & $T_e$ & Biagi v7.1 d+p group at 11.998 eV, as C1a\\
C1c & \el{} + Kr \ra \el{} + \Krs{} & $\langle\sigma v\rangle$, Table A.3 & $T_e$ & Biagi v7.1 higher group at 12.75 eV, as C1a\\
C2 & \el{} + Kr \ra 2\el{} + \Krp{} & $\langle\sigma v\rangle$, Table A.3 & $T_e$ & Biagi v7.1 ionization at 13.996 eV\\
G6 & \el{} + Kr, elastic & energy loss $(2m/M)\langle\sigma_m v\varepsilon\rangle$ & $T_e$, $T_g$ & Biagi v7.1 momentum transfer\\
ATT & \el{} + \Ft{} \ra F\ensuremath{^-} + F & $\langle\sigma v\rangle$, Table A.3 & $T_e$ & SIGLO F$_2$ set\\
DIS\_a3 & \el{} + \Ft{} \ra \el{} + 2F & $\langle\sigma v\rangle$, Table A.3 & $T_e$ & SIGLO {\unsansmath$a\,^3\Pi_u$}, repulsive, so excitation dissociates\\
DIS\_A1 & \el{} + \Ft{} \ra \el{} + 2F & $\langle\sigma v\rangle$, Table A.3 & $T_e$ & SIGLO {\unsansmath$A\,^1\Pi_u$}, as DIS\_a3\\
VIB1 & \el{} + \Ft{} \ra \el{} + \Ft{} & $\langle\sigma v\rangle$, Table A.3 & $T_e$ & SIGLO, v = 1\\
VIB2 & \el{} + \Ft{} \ra \el{} + \Ft{} & $\langle\sigma v\rangle$, Table A.3 & $T_e$ & SIGLO, v = 2\\
VIB3 & \el{} + \Ft{} \ra \el{} + \Ft{} & $\langle\sigma v\rangle$, Table A.3 & $T_e$ & SIGLO, v = 3\\
VIB4 & \el{} + \Ft{} \ra \el{} + \Ft{} & $\langle\sigma v\rangle$, Table A.3 & $T_e$ & SIGLO, v = 4\\
G2 & 2\Krs{} \ra \Krp{} + Kr + \el{} & $5.0{\times}10^{-10}$ cm$^3$s$^{-1}$, band kr\_pen &  & Kannari 1985 Table I Penning block, from Leichner and Ericson 1974. Band down to Werner's 1.5$\times10^{-10}$, via Jancaitis 1983\\
G3 & \Krs{} + \Ft{} \ra KrF(B) + F & $4.73{\times}10^{-10}$ cm$^3$s$^{-1}$ &  & rate from Velazco, Kolts and Setser 1976 and Ueda 1984, KrF$^*$ yield 0.90 from Velazco, Kolts and Setser 1976, B share from Yu and Setser 1990\\
G3C & \Krs{} + \Ft{} \ra KrF(C) + F & $2.9{\times}10^{-10}$ cm$^3$s$^{-1}$ &  & as G3, with Yu and Setser's $k_B/k_C$ of 1.6\\
G3X & \Krs{} + \Ft{} \ra Kr + 2F & $8.47{\times}10^{-11}$ cm$^3$s$^{-1}$ &  & as G3, the KrF$^*$-free rest, assigned to Kr + 2F\\
K15 & \Krp{} + 2Kr \ra \Krtp{} + Kr & $2.4{\times}10^{-31}$ cm$^6$s$^{-1}$ & $(T_g/300)^{-0.75}$ & the mean of three measurements, 2.4$\times10^{-31}$. $T_g$ law from Perez-Rios and Greene 2018\\
R11p & \el{} + \Krtp{} \ra \Krs{} + Kr & $1.6{\times}10^{-6}$ $(T_e/300\,\mathrm{K})^{-0.55}$ cm$^3$s$^{-1}$ & $T_e^{-0.55}$ & Shiu and Biondi 1977, 1.6$\times10^{-6}$ $(300\,\mathrm{K}/T_e)^{0.55}$, $\pm$ 10\,\%. Mikus 1978 within 8\,\%\\
K26 & \Krp{} + F\ensuremath{^-} \ra KrF(B) & $2.95{\times}10^{-6}$ cm$^3$s$^{-1}$ at 1\,atm, 300\,K, unscreened & $N$, screening, $(T_g/300)^{-3}$ on the 3-body part & Kannari 1985, Flannery turnover with Debye screening, and $(T_g/300)^{-3}$ on its three-body part from Thomson theory. Measured with F$^-$ only below 300 Torr, Jacob et al. 1978. No measurement near 1 atm\\
K108 & \Krtp{} + F\ensuremath{^-} \ra KrF(B) + Kr & $2.95{\times}10^{-6}$ cm$^3$s$^{-1}$ at 1\,atm, 300\,K, unscreened & $N$, screening, $(T_g/300)^{-3}$ on the 3-body part & as K26\\
G4 & KrF(B) \ra Kr + F + \hnu & $1.38{\times}10^{8}$ s$^{-1}$ &  & Kannari 1985 Table I, the 96.6\,\% B-X share of $\tau_B$ = 6.99 ns. Eden 1978 measured 6.8 $\pm$ 0.2 ns\\
G4A & KrF(B) \ra Kr + F + \hnu & $4.84{\times}10^{6}$ s$^{-1}$ &  & $A_{B\text{-}A}/A_{B\text{-}X}$ = 0.035, Dunning and Hay's calculation via Yu and Setser 1990. 272 nm is a theoretical position\\
G4C & KrF(C) \ra Kr + F + \hnu & $2.0{\times}10^{7}$ s$^{-1}$ &  & ML joint fit. Hakuta 1990 measured 60 $\pm$ 10 ns. Band at 275 nm, printed\\
BC & KrF(B) + Kr \ra KrF(C) + Kr & $1.05{\times}10^{-11}$ cm$^3$s$^{-1}$ &  & ML joint fit. Hakuta 1990's 1.2$\times10^{-10}$ is held out\\
CB & KrF(C) + Kr \ra KrF(B) + Kr & $9.55{\times}10^{-12}$ cm$^3$s$^{-1}$ & $1/K_\mathrm{eq}(T_g)$ & detailed balance of BC, with the joint fit's $K_\mathrm{eq}$ and Yu and Setser's measured 120 $\pm$ 50 cm$^{-1}$ gap\\
G5 & KrF(B) + 2Kr \ra \KtF{} + Kr & $1.9{\times}10^{-30}$ cm$^6$s$^{-1}$ & $(T_g/300)^{-3.32}$ & ML joint fit. $T_g$ law from Shui's calculation for 300-800 K, measured only to 425 K\\
G5C & KrF(C) + 2Kr \ra \KtF{} + Kr & $3.79{\times}10^{-31}$ cm$^6$s$^{-1}$ & $(T_g/300)^{-3.32}$ & ML joint fit\\
QE & KrF(B) + \el{} \ra Kr + F + \el{} & $6.6{\times}10^{-8}$ cm$^3$s$^{-1}$ &  & Trainor and Jacob's measured product, re-read with this network's lifetime\\
QEC & KrF(C) + \el{} \ra Kr + F + \el{} & $6.6{\times}10^{-8}$ cm$^3$s$^{-1}$ &  & as QE\\
QF2 & KrF(B) + \Ft{} \ra Kr + 3F & $4.75{\times}10^{-10}$ cm$^3$s$^{-1}$ &  & ML joint fit. Xu 1993 printed 3.0 $\pm$ 0.3$\times10^{-10}$ with fast mixing\\
QF2C & KrF(C) + \Ft{} \ra Kr + 3F & $3.04{\times}10^{-10}$ cm$^3$s$^{-1}$ &  & ML joint fit\\
QKR & KrF(B) + Kr \ra 2Kr + F & $1.0{\times}10^{-14}$ cm$^3$s$^{-1}$ &  & ML joint fit, at its lower bound\\
QKRC & KrF(C) + Kr \ra 2Kr + F & $8.37{\times}10^{-13}$ cm$^3$s$^{-1}$ &  & ML joint fit\\
C3 & \KtF{} \ra 2Kr + F + \hnu & $4.5{\times}10^{6}$ s$^{-1}$ &  & Xu, Gadomski and Setser 1993, 220 $\pm$ 20 ns. Band at 420 nm, printed\\
Q2F & \KtF{} + \Ft{} \ra 2Kr + F + \Ft{} & $1.2{\times}10^{-10}$ cm$^3$s$^{-1}$, band kr2f\_f2 &  & Xu 1993 Table I, 1.2 $\pm$ 0.2$\times10^{-10}$. Band up to Quigley and Hughes' 4.3$\times10^{-10}$\\
Q2FF & \KtF{} + F \ra 2Kr + 2F & 0, band kr2f\_f &  & no Kr measurement, so 0 unless a band run sets it, 0 to 1.8$\times10^{-10}$. The top is Jancaitis 1983's Ar$_2$F$^*$ analog\\
FF & 2F + Kr \ra \Ft{} + Kr & $2.8{\times}10^{-34}$ cm$^6$s$^{-1}$ &  & Cohen and Bott 1976 evaluation, 2.8$\times10^{-34}$, every rare gas given the Ar value. Held at its 300 K value\\
AK8\_Kr & \Krs{} + 2Kr \ra \Krts{} + Kr & $1.7{\times}10^{-32}$ cm$^6$s$^{-1}$, band kr2\_form & $(T_g/300)^{-0.64}$ & Ueda 1984 Table V, 1.7 $\pm$ 0.2$\times10^{-32}$ at 0.5-1 atm. $T_g$ law from Wieme and Lenaerts 1981, 300-451 K. Band up to Hughes 1977, 5.5$\times10^{-32}$\\
KR2\_RAD & \Krts{} \ra 2Kr + \hnu & $3.7{\times}10^{6}$ s$^{-1}$ &  & Morikawa 1989, 270 $\pm$ 5 ns, and five other measurements\\
NK7 & \Krts{} + \Ft{} \ra KrF(B) + Kr + F & $1.05{\times}10^{-10}$ cm$^3$s$^{-1}$, band kr2\_f2, half &  & half of Quigley and Hughes 1978, 2.1 $\pm$ 0.2$\times10^{-10}$. The even split is Kannari's. Band up to Chen and Payne's 5.0$\times10^{-10}$\\
K2F\_BR & \Krts{} + \Ft{} \ra \KtF{} + F & $1.05{\times}10^{-10}$ cm$^3$s$^{-1}$, band kr2\_f2, half &  & the other half of Quigley and Hughes 1978\\
K2F\_F & \Krts{} + F \ra KrF(B) + Kr & 0, band kr2\_f &  & no measurement, so 0 unless a band run sets it, 0 to 1$\times10^{-9}$. Kannari 1985's 3.0$\times10^{-10}$ is an estimate\\
PEN\_KR2 & 2\Krts{} \ra \Krtp{} + 2Kr + \el{} & 0, band kr2\_pen &  & no measurement, so 0 unless a band run sets it, 0 to 5.0$\times10^{-10}$. Kannari 1985 took 5.0$\times10^{-10}$ from an unpublished 1973 SRI report\\}
\def\wirGroups{KSTEP & \el{} + \Krs{} \ra \Krp{} + 2\el{} & $1.56{\times}10^{-7}\,T_e^{0.71}\times$ $e^{-4.09/T_e}$ cm$^3$s$^{-1}$ & $T_e$ & Chen R9 template, an estimate\\
SUPER & \el{} + \Krs{} \ra \el{} + Kr & from G1 by detailed balance & $T_e$ & detailed balance of G1, statistical weight ratio 1/12\\}
\def\wirVzero{G2 & 2\Krs{} \ra \Krp{} + Kr + \el{} & $5.0{\times}10^{-10}$ cm$^3$s$^{-1}$ &  & Kannari 1985 Table I, Penning block\\
G3 & \Krs{} + \Ft{} \ra KrF(B) + F & $8.1{\times}10^{-10}$ cm$^3$s$^{-1}$ &  & Kannari 1985 Table I, all to B. The report it credits gives 7.2$\times10^{-10}$\\
K15 & \Krp{} + 2Kr \ra \Krtp{} + Kr & $2.5{\times}10^{-31}$ cm$^6$s$^{-1}$ & $(T_g/300)^{-0.75}$ & Kannari 1985, Boichenko 2000\\
R11p & \el{} + \Krtp{} \ra \Krs{} + Kr & $2.24{\times}10^{-7}\,T_e^{-1/2}$ cm$^3$s$^{-1}$ & $T_e$ & Kannari 1985 Table I, 1.39$\times10^{-6}$ at 300 K on a $T_e^{-1/2}$ template\\
K26 & \Krp{} + F\ensuremath{^-} \ra KrF(B) & $2.95{\times}10^{-6}$ cm$^3$s$^{-1}$ at 1\,atm, 300\,K, unscreened & $N$, screening & Kannari 1985, Flannery turnover with Debye screening, no $T_g$ law\\
K108 & \Krtp{} + F\ensuremath{^-} \ra KrF(B) + Kr & $2.95{\times}10^{-6}$ cm$^3$s$^{-1}$ at 1\,atm, 300\,K, unscreened & $N$, screening & as K26\\
G4 & KrF(B) \ra Kr + F + \hnu & $1.43{\times}10^{8}$ s$^{-1}$ &  & Kannari 1985 Table I, all of $\tau_B$ as B-X\\
G5 & KrF(B) + 2Kr \ra \KtF{} + Kr & $6.0{\times}10^{-31}$ cm$^6$s$^{-1}$ & $(T_g/300)^{-3.32}$ & Xu, Gadomski and Setser 1993 Table II, 6.0 $\pm$ 2.0$\times10^{-31}$\\
G5\_C & KrF(C) + 2Kr \ra \KtF{} + Kr & $6.0{\times}10^{-31}$ cm$^6$s$^{-1}$ & $(T_g/300)^{-3.32}$ & as G5, applied to KrF(C)\\
AK8\_Kr & \Krs{} + 2Kr \ra \Krts{} + Kr & $2.7{\times}10^{-32}$ cm$^6$s$^{-1}$ &  & REBUILD's estimate, Kannari's Ar-partner rate given to Kr\\
KR2\_RAD & \Krts{} \ra 2Kr + \hnu & $8.0{\times}10^{7}$ s$^{-1}$ &  & Kannari 1985 Table I, 8.0$\times10^{7}$ per second\\
NK7 & \Krts{} + \Ft{} \ra KrF(B) + Kr + F & $3.0{\times}10^{-10}$ cm$^3$s$^{-1}$ &  & Kannari 1985 Table I\\
K2F\_BR & \Krts{} + \Ft{} \ra \KtF{} + F & $3.0{\times}10^{-10}$ cm$^3$s$^{-1}$ &  & Kannari 1985 Table I\\
K2F\_F & \Krts{} + F \ra KrF(B) + Kr & $3.0{\times}10^{-10}$ cm$^3$s$^{-1}$ &  & Kannari 1985, an estimate. No measurement exists\\
PEN\_KR2 & 2\Krts{} \ra \Krtp{} + 2Kr + \el{} & $5.0{\times}10^{-10}$ cm$^3$s$^{-1}$ &  & Kannari 1985 Table I, from an unpublished 1973 SRI report that gives Ar$_2^*$ the same value. No measurement held\\
K33 & KrF(B) + Kr \ra 2Kr + F & $4.0{\times}10^{-12}$ cm$^3$s$^{-1}$ &  & Xu, Gadomski and Setser 1993, 4.0 $\pm$ 0.5$\times10^{-12}$\\
K33\_C & KrF(C) + Kr \ra 2Kr + F & $4.0{\times}10^{-12}$ cm$^3$s$^{-1}$ &  & as K33, applied to KrF(C)\\
K32 & KrF(B) + \Ft{} \ra Kr + 3F & $3.0{\times}10^{-10}$ cm$^3$s$^{-1}$ &  & Xu, Gadomski and Setser 1993, 3.0 $\pm$ 0.3$\times10^{-10}$\\
K32\_C & KrF(C) + \Ft{} \ra Kr + 3F & $3.0{\times}10^{-10}$ cm$^3$s$^{-1}$ &  & as K32, applied to KrF(C)\\
K67 & KrF(B) + \el{} \ra Kr + F + \el{} & $4.5{\times}10^{-8}$ cm$^3$s$^{-1}$ &  & Kushner, Hanson and Schneider 1989\\
K67\_C & KrF(C) + \el{} \ra Kr + F + \el{} & $4.5{\times}10^{-8}$ cm$^3$s$^{-1}$ &  & as K67, applied to KrF(C)\\
K125 & \KtF{} + \Ft{} \ra 2Kr + F + \Ft{} & $1.5{\times}10^{-10}$ cm$^3$s$^{-1}$ &  & Kannari 1985 Table I\\
K\_2Fe & \KtF{} + \el{} \ra 2Kr + F + \el{} & $1.0{\times}10^{-7}$ cm$^3$s$^{-1}$ &  & Kannari 1985, estimated there\\
BC\_Kr & KrF(B) + Kr \ra KrF(C) + Kr & $5.0{\times}10^{-10}$ cm$^3$s$^{-1}$ &  & Kannari's B-C mixing with a Kr partner, an estimate\\
CB\_Kr & KrF(C) + Kr \ra KrF(B) + Kr & $3.41{\times}10^{-10}$ cm$^3$s$^{-1}$ &  & detailed balance of BC\_Kr at 300 K\\
URAD\_C & KrF(C) \ra Kr + F + \hnu & $1.33{\times}10^{7}$ s$^{-1}$ &  & Kannari 1985 Table I, 1.33$\times10^{7}$ per second\\}
\def\wirXsec{G1 & Kr, 5s manifold, 9.915\,eV & $1.24{\times}10^{-17}$ & $3.44{\times}10^{-13}$ & $1.11{\times}10^{-11}$ & $6.41{\times}10^{-11}$ & $3.78{\times}10^{-10}$ & 0.036\\
C1a & Kr, 5p group, 11.3\,eV & $5.1{\times}10^{-19}$ & $5.93{\times}10^{-14}$ & $2.97{\times}10^{-12}$ & $2.11{\times}10^{-11}$ & $1.48{\times}10^{-10}$ & $5.53{\times}10^{-3}$\\
C1b & Kr, d+p group, 11.998\,eV & $2.5{\times}10^{-19}$ & $5.84{\times}10^{-14}$ & $3.77{\times}10^{-12}$ & $3.08{\times}10^{-11}$ & $2.53{\times}10^{-10}$ & $1.81{\times}10^{-3}$\\
C1c & Kr, higher group, 12.75\,eV & $9.15{\times}10^{-21}$ & $7.73{\times}10^{-15}$ & $8.8{\times}10^{-13}$ & $1.02{\times}10^{-11}$ & $1.29{\times}10^{-10}$ & $1.48{\times}10^{-4}$\\
C2 & Kr, ionization, 13.996\,eV & $8.06{\times}10^{-21}$ & $1.57{\times}10^{-14}$ & $2.28{\times}10^{-12}$ & $2.94{\times}10^{-11}$ & $4.14{\times}10^{-10}$ & $5.09{\times}10^{-5}$\\
ATT & F$_2$, attachment & $6.58{\times}10^{-9}$ & $3.51{\times}10^{-9}$ & $2.25{\times}10^{-9}$ & $1.61{\times}10^{-9}$ & $9.75{\times}10^{-10}$ & 1\\
DIS\_a3 & F$_2$, {\unsansmath$a\,^3\Pi_u$}, 3.16\,eV & $3.01{\times}10^{-12}$ & $1.13{\times}10^{-10}$ & $3.95{\times}10^{-10}$ & $7.37{\times}10^{-10}$ & $1.35{\times}10^{-9}$ & 1\\
DIS\_A1 & F$_2$, {\unsansmath$A\,^1\Pi_u$}, 4.34\,eV & $2.06{\times}10^{-13}$ & $1.93{\times}10^{-11}$ & $8.88{\times}10^{-11}$ & $1.9{\times}10^{-10}$ & $4.02{\times}10^{-10}$ & 1\\
VIB1 & F$_2$, $v=1$, 0.111\,eV & $1.77{\times}10^{-9}$ & $1.22{\times}10^{-9}$ & $8.72{\times}10^{-10}$ & $6.61{\times}10^{-10}$ & $4.27{\times}10^{-10}$ & 1\\
VIB2 & F$_2$, $v=2$, 0.219\,eV & $2.01{\times}10^{-9}$ & $1.49{\times}10^{-9}$ & $1.09{\times}10^{-9}$ & $8.37{\times}10^{-10}$ & $5.47{\times}10^{-10}$ & 1\\
VIB3 & F$_2$, $v=3$, 0.324\,eV & $1.21{\times}10^{-9}$ & $9.75{\times}10^{-10}$ & $7.35{\times}10^{-10}$ & $5.71{\times}10^{-10}$ & $3.78{\times}10^{-10}$ & 1\\
VIB4 & F$_2$, $v=4$, 0.421\,eV & $3.83{\times}10^{-10}$ & $3.27{\times}10^{-10}$ & $2.51{\times}10^{-10}$ & $1.96{\times}10^{-10}$ & $1.31{\times}10^{-10}$ & 1\\
\midrule
ATT & energy an attached \el{} takes, eV & 0.37 & 0.46 & 0.53 & 0.59 & 0.72 & \\
 & $\tfrac32 T_e$, for comparison & 0.75 & 1.50 & 2.25 & 3.00 & 4.50 & \\}
\def\wirTe{0.5 & 1 & 1.5 & 2 & 3}
\def\wirBands{kr2\_f & \Krts{} + F \ra KrF(B) + Kr & 0 & 0 to $1.0{\times}10^{-9}$ cm$^3$s$^{-1}$ & No measurement. Kannari's 3.0$\times10^{-10}$ traces to SRI MP-76-99, whose model leaves the reaction out (p.45) and whose one number is a conditional 3-10$\times10^{-10}$ for Ar$_2^*$ + F (p.65), the band's top. Jancaitis 1983 has no such row\\
trimer\_tg\_pow & the trimer's $(T_g/300)^n$ & $-3.32$ & $-3.32$ to $-0.50$ exponent & Shui's calculation over 300-800 K, the default, measured to 425 K. Salesky and Kimura inferred \ensuremath{-}0.5 from pure-Kr efficiencies, a bound that never sets it. Params.trimer\_tg\_pow sets it\\
kr\_pen & \Krs{} + \Krs{} \ra \Krp{} + Kr + \el{} & $5.0{\times}10^{-10}$ & $1.5{\times}10^{-10}$ to $5.0{\times}10^{-10}$ cm$^3$s$^{-1}$ & Kannari 1985 gives 5.0$\times10^{-10}$ from Leichner and Ericson 1974, v0's and the default. Jancaitis's model carries 1.5$\times10^{-10}$ from Werner's thesis (1975)\\
kr2f\_f2 & \KtF{} + \Ft{} & $1.2{\times}10^{-10}$ & $1.2{\times}10^{-10}$ to $4.3{\times}10^{-10}$ cm$^3$s$^{-1}$ & Xu 1993 measured 1.2 $\pm$ 0.2$\times10^{-10}$, the default, and lists 1.5, 1.8 and 4.3$\times10^{-10}$ before it. Jancaitis 1983 measured 1.8 $\pm$ 0.2$\times10^{-10}$\\
kr2f\_f & \KtF{} + F \ra 2Kr + 2F & 0 & 0 to $1.8{\times}10^{-10}$ cm$^3$s$^{-1}$ & No Kr measurement. Jancaitis 1983 found F atoms quench Ar$_2$F$^*$ at 0.5-1 times F$_2$ (p.153). With his Kr$_2$F$^*$ + F$_2$, 1.8$\times10^{-10}$, that analog gives 0.9-1.8$\times10^{-10}$ and sets the top\\
kr2\_form & \Krs{} + 2Kr \ra \Krts{} + Kr, 300\unit{K} & $1.7{\times}10^{-32}$ & $1.7{\times}10^{-32}$ to $5.5{\times}10^{-32}$ cm$^6$s$^{-1}$ & Ueda 1984 measured 1.7 $\pm$ 0.2$\times10^{-32}$ at 0.5-1 atm, the default. Low-pressure afterglows give 1.79 to 4.4$\times10^{-32}$ and Hughes 1977 5.5$\times10^{-32}$\\
kr2\_f2 & \Krts{} + \Ft{}, both channels & $2.1{\times}10^{-10}$ & $2.1{\times}10^{-10}$ to $5.0{\times}10^{-10}$ cm$^3$s$^{-1}$ & Quigley and Hughes 1978 measured 2.1 $\pm$ 0.2$\times10^{-10}$, the default. SRI MP-76-99 gives 3.0$\times10^{-10}$ and Chen and Payne 1979 5.0$\times10^{-10}$, beyond each other's errors\\
kr2\_pen & \Krts{} + \Krts{} \ra \Krtp{} + 2Kr + \el{} & 0 & 0 to $5.0{\times}10^{-10}$ cm$^3$s$^{-1}$ & No measurement. Kannari's 5.0$\times10^{-10}$ is Lorents and Olson's unpublished 1973 SRI report, which gives Ar$_2^*$ the same value, and Jancaitis's model carries 1.0$\times10^{-10}$ from Werner et al. 1977\\}
\def\wirSpecies{Kr & 0 & 1.5 & ground state\\
\Ft{} & 0 & 2.5 & ground state\\
F & 0.8 & 1.5 & half the F$_2$ bond, $D_0$ = 1.60 eV\\
\Krs{} & 9.915 & 1.5 & Biagi's 5s threshold, the 5s[3/2]2 level\\
\Krp{} & 13.996 & 1.5 & Biagi's ionization threshold\\
KrF(B) & 5.856 & 2.5 & $D_e(B)$ = 5.54 eV below the ion pair, Tellinghuisen 1976\\
KrF(C) & 5.84112 & 2.5 & 120 cm$^{-1}$ below B, Yu and Setser 1990\\
\KtF{} & 5.276 & 3 & 0.58 eV below KrF$^*$ + Kr, calculated by Huestis et al. 1984\\
\Krts{} & 9.2 & 2.5 & the 147-nm excimer, an estimated well\\
\Krtp{} & 12.846 & 2.5 & IP(Kr) less $D_0$(Kr$_2^+$) = 1.15 eV, measured by photoionization, Ng et al. 1977\\
F\ensuremath{^-} & \ensuremath{-}2.6 & 1.5 & $E_\mathrm{form}$(F) less the electron affinity, 3.40 eV\\}

% LIGHT v0_2 accuracy numbers (fast). GENERATED by make_overview_v0_2.py
% from LIGHT v0_2 (0.1.2) and ML/. Do not edit. Rerun the script instead.
\providecommand\widef[2]{\expandafter\def\csname wi@#1\endcsname{#2}}
\widef{xu.50}{0.91}
\widef{xu.50.v0}{2.89}
\widef{xu.100}{0.97}
\widef{xu.100.v0}{2.41}
\widef{xu.50.meas}{$5.0{\times}10^{7}$}
\widef{xu.100.meas}{$6.2{\times}10^{7}$}
\widef{eq5.dev}{$1.7{\times}10^{-4}$}
\widef{gap}{120}
\widef{closure.max}{$10^{-12}$}
\widef{closure.runs}{9}
\widef{hot.p}{0.1}
\widef{hot.P}{500}
\widef{hot.tg}{3458}
\widef{sc.eden_kr}{2.3}
\widef{sc.eden_f2}{3.2}
\widef{sc.ksu_fig2}{4.2}
\widef{sc.ksu_fig1}{6.1}
\widef{sc.xu_figs5_7}{11.7}
\widef{sc.hakuta_fig3}{34.1}
\widef{sc.hakuta_fig2}{19.3}
\widef{sc.ys_ir300}{\ensuremath{-}0.5}
\widef{sc.ys_ratio_230_300}{\ensuremath{-}8.4}
\widef{sc.bs_tau_b}{\ensuremath{-}22.3}
\widef{ys.meas}{0.16}
\widef{ys.err}{0.02}
\widef{ysr.meas}{1.30}
\widef{ysr.err}{0.06}
\widef{sc.train.lo}{2.3}
\widef{sc.train.hi}{6.1}
\widef{sc.bs.meas}{9.0}
\widef{sc.bs.err}{0.6}
\widef{tah.g0.meas}{6.39}
\widef{tah.g0}{5.69}
\widef{tah.g0.x}{0.89}
\widef{tah.isat.meas}{1.4}
\widef{tah.isat}{1.55}
\widef{tah.isat.err}{0.1}
\widef{tah.isat.x}{1.11}
\widef{tah.eta.meas}{0.27}
\widef{tah.eta}{0.268}
\widef{tah.eta.x}{0.99}
\widef{tah.form}{0.300}
\widef{tah.tg}{446}
\widef{tah.pin}{330}
\widef{tah.torr}{665}
\widef{tah.f2torr}{2.8}
\widef{tah.seed}{$1.0{\times}10^{10}$}
\widef{tah.def.g0}{5.69}
\widef{tah.win.lo}{200}
\widef{tah.win.hi}{400}
\widef{tah.f2left}{29}
\widef{tah.bd.kr2_f.lo}{5.69}
\widef{tah.bd.kr2_f.hi}{6.23}
\widef{tah.bd.trimer_tg_pow.lo}{5.69}
\widef{tah.bd.trimer_tg_pow.hi}{2.84}
\widef{tah.bd.kr_pen.lo}{5.71}
\widef{tah.bd.kr_pen.hi}{5.69}
\widef{tah.bd.kr2f_f2.lo}{5.69}
\widef{tah.bd.kr2f_f2.hi}{5.89}
\widef{tah.bd.kr2f_f.lo}{5.69}
\widef{tah.bd.kr2f_f.hi}{5.86}
\widef{tah.bd.kr2_form.lo}{5.69}
\widef{tah.bd.kr2_form.hi}{5.31}
\widef{tah.bd.kr2_f2.lo}{5.69}
\widef{tah.bd.kr2_f2.hi}{5.74}
\widef{tah.bd.kr2_pen.lo}{5.69}
\widef{tah.bd.kr2_pen.hi}{5.69}
\widef{tah.bd.minor}{0.38}
\widef{tah.bd.widest}{the trimer's $(T_g/300)^n$}
\widef{tah.hiF.g0}{6.23}
\widef{tah.hiF.x}{0.97}
\widef{tah.tri.g0}{2.84}
\widef{tah.tri.x}{0.44}
\widef{tri.lo.pow}{0.5}
\widef{tri.klimek.hi}{425}
\widef{tah.kannari.g0}{6.10}
\widef{tah.kannari.x}{0.95}
\widef{tah.own.g0}{5.53}
\widef{tah.own.te}{3.8}
\widef{tah.own.ne}{$2.59{\times}10^{14}$}
\widef{tah.own.tg}{423}
\widef{tah.beam.g0}{5.69}
\widef{tah.beam.te}{1.6}
\widef{tah.beam.ne}{$6.33{\times}10^{14}$}
\widef{tah.beam.tg}{446}
\widef{tah.own.x}{0.87}
\widef{kr2.def}{4.67}
\widef{kr2.vuv}{20.0}
\widef{kr2.burn}{640}
\widef{kr2.F}{$2.12{\times}10^{17}$}
\widef{kr2.lo.k}{$1.0{\times}10^{-11}$}
\widef{kr2.lo}{6.52}
\widef{kr2.kann.k}{$3.0{\times}10^{-10}$}
\widef{kr2.kann}{9.21}
\widef{kr2.kann.ratio}{17}
\widef{kr2.hi.k}{$1.0{\times}10^{-9}$}
\widef{kr2.hi}{9.50}
\widef{kr2.hi.vuv}{0.5}
\widef{kr2.hi.ratio}{57}
\widef{kr2.hi.afterpc}{74}
\widef{kr2.both}{9.17}
\widef{kr2.pen.k}{$5.0{\times}10^{-10}$}
\widef{pure.Kr.x}{1.02}
\widef{pure.Ar.x}{1.03}
\widef{pure.dr.lo}{1.00}
\widef{pure.dr.hi}{1.00}
\widef{pure.dr.te.lo}{0.026}
\widef{pure.dr.te.hi}{2}
\widef{pure.k15}{$2.4{\times}10^{-31}$}
\widef{pure.conv.lo}{$2.74{\times}10^{-31}$}
\widef{pure.conv.hi}{$3.3{\times}10^{-31}$}
\widef{pure.vuv}{0.57}
\widef{pure.vuv.meas}{0.42}
\widef{pure.vuv.err}{0.05}
\widef{pure.vuv.x}{1.36}
\widef{pure.vuv.ideal}{0.54}
\widef{pure.vuv.nothermal}{0.55}
\widef{pure.thermal}{0.060}
\widef{pure.band.move}{0.001}
\widef{pure.hi.P}{300}
\widef{pure.hi.pulse}{1}
\widef{pure.hi.torr}{665}
\widef{pure.hi.def}{0.27}
\widef{pure.hi.lo}{0.13}
\widef{pure.hi.lo.what}{\Krts{} + \Krts{} \ra \Krtp{} + 2Kr + \el{} at $5.0{\times}10^{-10}$\unit{cm^3\,s^{-1}}}
\widef{pure.hi.hi}{0.49}
\widef{pure.hi.hi.what}{\Krs{} + 2Kr \ra \Krts{} + Kr at $5.5{\times}10^{-32}$\unit{cm^6\,s^{-1}}}
\def\wiTahBands{\Krts{} + F \ra KrF(B) + Kr & 5.69 to 6.23 & \x{0.89} to \x{0.97}\\
the trimer's $(T_g/300)^n$ & 5.69 to 2.84 & \x{0.89} to \x{0.44}\\
\Krs{} + \Krs{} \ra \Krp{} + Kr + \el{} & 5.71 to 5.69 & \x{0.89} to \x{0.89}\\
\KtF{} + \Ft{} & 5.69 to 5.89 & \x{0.89} to \x{0.92}\\
\KtF{} + F \ra 2Kr + 2F & 5.69 to 5.86 & \x{0.89} to \x{0.92}\\
\Krs{} + 2Kr \ra \Krts{} + Kr, 300\unit{K} & 5.69 to 5.31 & \x{0.89} to \x{0.83}\\
\Krts{} + \Ft{}, both channels & 5.69 to 5.74 & \x{0.89} to \x{0.90}\\
\Krts{} + \Krts{} \ra \Krtp{} + 2Kr + \el{} & 5.69 to 5.69 & \x{0.89} to \x{0.89}\\}

% LIGHT v0_2 accuracy numbers (slow). GENERATED by make_overview_v0_2.py
% from LIGHT v0_2 (0.1.2) and ML/. Do not edit. Rerun the script instead.
\providecommand\widef[2]{\expandafter\def\csname wi@#1\endcsname{#2}}
\widef{fit.k_BC.lo}{$4.63{\times}10^{-12}$}
\widef{fit.k_BC.mid}{$9.8{\times}10^{-12}$}
\widef{fit.k_BC.hi}{$1.73{\times}10^{-11}$}
\widef{fit.K_eq.lo}{0.84}
\widef{fit.K_eq.mid}{1.07}
\widef{fit.K_eq.hi}{1.38}
\widef{fit.tau_C.lo}{41}
\widef{fit.tau_C.mid}{50}
\widef{fit.tau_C.hi}{63}
\widef{fit.tau_B.lo}{6.79}
\widef{fit.tau_B.mid}{7.02}
\widef{fit.tau_B.hi}{7.22}
\widef{fit.k_F2.lo}{$3.92{\times}10^{-10}$}
\widef{fit.k_F2.mid}{$4.57{\times}10^{-10}$}
\widef{fit.k_F2.hi}{$5.3{\times}10^{-10}$}
\widef{fit.k_F2_C.lo}{$2.74{\times}10^{-10}$}
\widef{fit.k_F2_C.mid}{$3.02{\times}10^{-10}$}
\widef{fit.k_F2_C.hi}{$3.32{\times}10^{-10}$}
\widef{fit.k3_Kr.lo}{$1.26{\times}10^{-30}$}
\widef{fit.k3_Kr.mid}{$1.92{\times}10^{-30}$}
\widef{fit.k3_Kr.hi}{$2.44{\times}10^{-30}$}
\widef{fit.k3_Kr_C.lo}{$1.93{\times}10^{-31}$}
\widef{fit.k3_Kr_C.mid}{$3.64{\times}10^{-31}$}
\widef{fit.k3_Kr_C.hi}{$5.12{\times}10^{-31}$}
\widef{fit.bshare.lo}{0.57}
\widef{fit.bshare.mid}{0.69}
\widef{fit.bshare.hi}{0.80}
\widef{fit.bshare.light}{0.68}
\widef{fit.n}{10}
\widef{post.g0.lo}{4.99}
\widef{post.g0.mid}{5.39}
\widef{post.g0.hi}{6.18}
\widef{post.n}{40}
\widef{link.base}{0.268}
\widef{link.yield.lo}{0.240}
\widef{link.yield.lo.d}{\ensuremath{-}0.028}
\widef{link.yield.hi}{0.282}
\widef{link.yield.hi.d}{+0.014}
\widef{link.yield.arf}{0.211}
\widef{link.yield.arf.d}{\ensuremath{-}0.057}
\widef{link.nascentC}{0.290}
\widef{link.nascentC.d}{+0.022}
\widef{link.ionY}{0.225}
\widef{link.ionY.d}{\ensuremath{-}0.043}
\widef{link.ionC}{0.251}
\widef{link.ionC.d}{\ensuremath{-}0.017}
\widef{link.ionRate}{0.266}
\widef{link.ionRate.d}{\ensuremath{-}0.002}
\widef{link.kr2F.kann}{0.285}
\widef{link.kr2F.kann.d}{+0.017}
\widef{link.kr2F.hi}{0.290}
\widef{link.kr2F.hi.d}{+0.022}
\widef{link.kr2split.lo}{0.260}
\widef{link.kr2split.lo.d}{\ensuremath{-}0.008}
\widef{link.kr2split.hi}{0.276}
\widef{link.kr2split.hi.d}{+0.008}
\widef{link.k3}{0.256}
\widef{link.k3.d}{\ensuremath{-}0.012}
\widef{link.minor}{0.013}
\widef{el.loop.te}{1.61}
\widef{el.loop.ne}{$6.49{\times}10^{14}$}
\widef{el.loop.g0}{5.63}
\widef{el.loop.form}{0.299}
\widef{el.light.te}{1.60}
\widef{el.light.ne}{$6.33{\times}10^{14}$}
\widef{el.light.g0}{5.69}
\widef{el.light.form}{0.300}
\widef{el.dev.te}{0.7}
\widef{el.dev.ne}{2.5}
\widef{el.dev.g0}{1.0}
\widef{el.dev.formation}{0.5}
\widef{el.eatt}{0.55}
\widef{el.eatt.32}{2.41}
\widef{el.rpow}{0.55}
\widef{el.rcap}{0.95}
\widef{sw.krw.lo}{0.96}
\widef{sw.krw.hi}{1.03}
\widef{sw.f2.lo}{0.48}
\widef{sw.f2.hi}{2.41}
\def\wiLinks{KrF$^*$ yield of \Krs{} + \Ft & 0.90 & 0.70 & the measured 0.90 less its stated 0.20 & 0.240 & \ensuremath{-}0.028\\
 &  & 1.0 & the published assignment & 0.282 & +0.014\\
 &  & 0.5 & as for Ar$^*$ + F$_2$ & 0.211 & \ensuremath{-}0.057\\
share of that KrF formed in C & 0.38 & 0 & if nascent C relaxes to B at 1 atm & 0.290 & +0.022\\
KrF$^*$ yield of ion-ion recombination & 1 & 0.64 & as for Cl$^-$ & 0.225 & \ensuremath{-}0.043\\
share of that KrF formed in C & 0 & 0.38 & as for SF$_6^-$ & 0.251 & \ensuremath{-}0.017\\
ion-ion recombination rate & Flannery & \x{0.38} & Avco's model input & 0.266 & \ensuremath{-}0.002\\
\Krts{} + F \ra KrF(B) + Kr & 0 & $3.0{\times}10^{-10}$ & Kannari's estimate & 0.285 & +0.017\\
 &  & $1.0{\times}10^{-9}$ & the top of its band & 0.290 & +0.022\\}

% LIGHT v0_2 energy ladder. GENERATED by make_energy_ladder.py
% from LIGHT v0_2 (0.1.2). Do not edit. Rerun the script instead.
\providecommand\laddef[2]{\expandafter\def\csname lad@#1\endcsname{#2}}
\laddef{version}{0.1.2}
\laddef{E.B}{5.856}
\laddef{DeB}{5.54}
\laddef{BoverKrF}{5.06}
\laddef{wl.G4}{248}
\laddef{hv.G4}{5.00}
\laddef{heat.G4}{0.06}
\laddef{heat.G3}{3.26}
\laddef{heat.K26}{5.54}
\laddef{heat.G5}{0.58}
\laddef{wl.C3}{420}
\laddef{hv.C3}{2.95}
\laddef{heat.C3}{1.53}
\laddef{gapBC}{120}
\laddef{notdrawn}{G6, VIB1, VIB2, VIB3 and VIB4}
\laddef{toe}{G2, PEN\_KR2 and SUPER}
\laddef{zero}{Q2FF, K2F\_F and PEN\_KR2}
\laddef{groups}{\texttt{electron}}
\laddef{pale}{}
\laddef{wl.V}{147}
\laddef{hv.V}{8.43}
\laddef{heat.V}{0.77}
\def\ladpicture{%
\definecolor{ladE}{HTML}{2A78D6}% electron gas
\definecolor{ladX}{HTML}{EDA100}% Kr*
\definecolor{ladB}{HTML}{1BAF7A}% KrF(B), 248 and 272 nm
\definecolor{ladC}{HTML}{008300}% KrF(C), 275 nm
\definecolor{ladT}{HTML}{4A3AA7}% Kr2F*, 420 nm
\definecolor{ladV}{HTML}{A63A83}% Kr2*, 147 nm
\definecolor{ladH}{HTML}{E34948}% heat
\definecolor{ladI}{HTML}{52514E}% ions
\definecolor{ladF}{HTML}{A8A69E}% F atoms
\definecolor{ladInk}{HTML}{0B0B0B}%
\definecolor{ladInk2}{HTML}{52514E}%
\definecolor{ladMuted}{HTML}{898781}%
\tikzset{
  lad grid/.style={draw=ladMuted!22, line width=0.3pt},
  lad axisline/.style={draw=ladMuted, line width=0.5pt},
  lad lev/.style={line width=1.4pt},
  lad pair/.style={line width=0.8pt, dash pattern=on 3pt off 2pt},
  lad tick/.style={line width=0.5pt},
  lad arr/.style={line width=0.7pt, -{Triangle[width=3.2pt,length=3.4pt]}},
  lad hv/.style={line width=1.3pt, -{Triangle[width=4.2pt,length=4.4pt]}},
  lad on/.style={},
  lad opt/.style={dash pattern=on 2.4pt off 1.6pt},
  lad zero/.style={dash pattern=on 0.6pt off 1.4pt, line cap=round},
  lad heat/.style={draw=ladH, line width=0.6pt, decorate,
    decoration={snake, amplitude=0.5mm, segment length=1.4mm,
    post length=0.5mm}},
  lad lab/.style={font=\scriptsize, text=ladInk, inner sep=1pt},
  lad lab2/.style={font=\scriptsize, text=ladInk, inner sep=1pt,
    align=right},
  lad lab2l/.style={font=\scriptsize, text=ladInk, inner sep=1pt,
    align=left},
  lad alab/.style={font=\scriptsize, text=ladInk, inner sep=0.8pt,
    fill=white, fill opacity=0.85, text opacity=1, align=center},
  lad slab/.style={font=\tiny, text=ladInk, inner sep=0.6pt,
    fill=white, fill opacity=0.85, text opacity=1, align=right},
  lad hlab/.style={font=\tiny, text=ladInk2, inner sep=0.8pt,
    fill=white, fill opacity=0.85, text opacity=1},
  lad tlab/.style={font=\tiny, text=ladInk2, inner sep=0.8pt},
  lad axis/.style={font=\scriptsize, text=ladInk2, inner sep=1pt},
  lad head/.style={font=\footnotesize\bfseries, text=ladInk2},
}
\begin{tikzpicture}
\draw[lad grid] (0.950,1.000) -- (16.200,1.000);
\draw[lad grid] (0.950,3.000) -- (16.200,3.000);
\draw[lad grid] (0.950,5.000) -- (16.200,5.000);
\draw[lad grid] (0.950,7.000) -- (16.200,7.000);
\draw[lad grid] (0.950,9.000) -- (16.200,9.000);
\draw[lad grid] (0.950,11.000) -- (16.200,11.000);
\draw[lad grid] (0.950,13.000) -- (16.200,13.000);
\draw[lad grid] (0.950,15.000) -- (16.200,15.000);
\draw[lad grid] (0.950,17.000) -- (16.200,17.000);
\draw[lad axisline] (0.950,0.000) -- (0.950,17.450);
\draw[lad lev,draw=ladInk2] (2.650,3.000) -- (6.300,3.000);
\draw[lad lev,draw=ladX] (2.650,12.915) -- (6.300,12.915);
\draw[lad lev,draw=ladI] (2.650,16.996) -- (6.300,16.996);
\draw[lad lev,draw=ladV] (5.000,12.200) -- (6.300,12.200);
\draw[lad lev,draw=ladI] (5.000,15.846) -- (6.300,15.846);
\draw[lad tick,draw=ladMuted] (3.550,14.300) -- (4.400,14.300);
\draw[lad tick,draw=ladMuted] (3.730,14.998) -- (4.400,14.998);
\draw[lad tick,draw=ladMuted] (3.910,15.750) -- (4.400,15.750);
\draw[lad tick,draw=ladMuted] (5.740,3.770) -- (6.020,3.770);
\draw[lad pair,draw=ladInk2] (10.200,14.396) -- (12.300,14.396);
\draw[lad pair,draw=ladInk2] (10.200,13.246) -- (12.300,13.246);
\draw[lad pair,draw=ladInk2] (7.000,3.800) -- (12.300,3.800);
\draw[lad lev,draw=ladB] (8.400,8.886) -- (11.050,8.886);
\draw[lad lev,draw=ladC] (8.400,8.811) -- (11.050,8.811);
\draw[lad lev,draw=ladT] (7.300,8.276) -- (8.650,8.276);
\draw[lad tick,draw=ladMuted] (9.210,3.856) -- (9.490,3.856);
\draw[lad tick,draw=ladMuted] (9.610,4.296) -- (9.890,4.296);
\draw[lad tick,draw=ladMuted] (10.010,4.331) -- (10.290,4.331);
\draw[lad tick,draw=ladMuted] (7.840,5.326) -- (8.120,5.326);
\draw[lad lev,draw=ladInk2] (13.900,3.000) -- (15.100,3.000);
\draw[lad lev,draw=ladF] (13.900,3.800) -- (15.100,3.800);
\draw[lad pair,draw=ladInk2] (13.900,4.600) -- (15.100,4.600);
\draw[lad pair,draw=ladInk2] (13.900,1.200) -- (15.100,1.200);
\draw[lad lev,draw=ladI] (13.900,0.400) -- (15.100,0.400);
\draw[lad tick,draw=ladMuted] (14.290,6.160) -- (15.100,6.160);
\draw[lad tick,draw=ladMuted] (14.570,7.340) -- (15.100,7.340);
\draw[lad arr,lad on,draw=ladE] (2.800,3.000) -- node[pos=0.52,rotate=90,lad alab] {C2} (2.800,16.996);
\draw[lad arr,lad on,draw=ladE] (3.080,3.000) -- node[pos=0.40,rotate=90,lad alab] {G1} (3.080,12.915);
\draw[lad arr,lad opt,draw=ladE] (3.360,12.915) -- node[pos=0.78,rotate=90,lad alab] {SUPER} (3.360,3.000);
\draw[lad arr,lad on,draw=ladX] (3.080,12.915) -- node[pos=0.36,rotate=90,lad alab] {G2} (3.080,16.996);
\draw[lad arr,lad opt,draw=ladE] (3.360,12.915) -- node[pos=0.62,rotate=90,lad alab] {KSTEP} (3.360,16.996);
\draw[lad arr,lad on,draw=ladE] (3.620,3.000) -- (3.620,14.300);
\draw[lad arr,lad on,draw=ladE] (3.800,3.000) -- node[pos=0.55,rotate=90,lad alab] {C1a--c} (3.800,14.998);
\draw[lad arr,lad on,draw=ladE] (3.980,3.000) -- (3.980,15.750);
\draw[lad heat] (4.400,15.750) -- (4.400,12.915);
\draw[lad arr,lad on,draw=ladH] (2.700,12.915) -- (2.350,12.065);
\draw[lad arr,lad on,draw=ladX] (5.100,12.915) -- (5.100,12.200);
\draw[lad arr,lad on,draw=ladI] (5.250,16.996) -- node[pos=0.50,rotate=90,lad alab] {K15} (5.250,15.846);
\draw[lad arr,lad on,draw=ladI] (5.800,15.846) -- node[pos=0.45,rotate=90,lad alab] {R11p} (5.800,12.915);
\draw[lad arr,lad zero,draw=ladV] (6.120,12.200) -- node[pos=0.62,rotate=90,lad alab] {PEN\_KR2} (6.120,15.846);
\draw[lad hv,lad on,draw=ladV] (5.880,12.200) -- node[pos=0.50,rotate=90,lad alab] {KR2\_RAD, 147 nm} (5.880,3.770);
\draw[lad heat] (5.880,3.770) -- (5.880,3.000);
\draw[lad arr,lad on,draw=ladV] (6.300,12.200) -- node[pos=0.42,sloped,lad alab] {K2F\_BR} (7.620,8.276);
\draw[lad arr,lad on,draw=ladV] (6.300,12.200) -- node[pos=0.42,sloped,lad alab] {NK7} (8.660,8.856);
\draw[lad arr,lad zero,draw=ladV] (6.300,12.200) -- node[pos=0.62,sloped,lad alab] {K2F\_F} (8.900,8.856);
\draw[lad arr,lad on,draw=ladX] (6.300,12.915) -- node[pos=0.40,sloped,lad alab] {G3, G3C} (9.140,8.856);
\draw[lad arr,lad on,draw=ladI] (10.450,13.246) -- node[pos=0.50,rotate=90,lad alab] {K108} (10.450,8.856);
\draw[lad arr,lad on,draw=ladI] (10.720,14.396) -- node[pos=0.62,rotate=90,lad alab] {K26} (10.720,8.856);
\draw[lad arr,lad on,draw=ladB] (8.520,8.856) -- (8.520,8.276);
\draw[lad hv,lad on,draw=ladB] (9.350,8.856) -- node[pos=0.50,rotate=90,lad alab] {G4, 248 nm} (9.350,3.856);
\draw[lad hv,lad on,draw=ladB] (9.750,8.856) -- node[pos=0.50,rotate=90,lad alab] {G4A, 272 nm} (9.750,4.296);
\draw[lad heat] (9.750,4.296) -- (9.750,3.800);
\draw[lad hv,lad on,draw=ladC] (10.150,8.841) -- node[pos=0.50,rotate=90,lad alab] {G4C, 275 nm} (10.150,4.331);
\draw[lad heat] (10.150,4.331) -- (10.150,3.800);
\draw[lad arr,lad on,draw=ladH] (10.930,8.841) -- (10.930,8.041);
\draw[lad hv,lad on,draw=ladT] (7.980,8.276) -- node[pos=0.64,rotate=90,lad alab] {C3, 420 nm} (7.980,5.326);
\draw[lad heat] (7.980,5.326) -- (7.980,3.800);
\draw[lad arr,lad on,draw=ladH] (7.520,8.276) -- (7.520,7.476);
\draw[lad arr,lad zero,draw=ladH] (8.400,8.276) -- (8.400,7.476);
\draw[lad arr,lad on,draw=ladE] (14.360,3.000) -- node[pos=0.62,rotate=90,lad alab] {DIS\_a3} (14.360,6.160);
\draw[lad arr,lad on,draw=ladE] (14.640,3.000) -- node[pos=0.66,rotate=90,lad alab] {DIS\_A1} (14.640,7.340);
\draw[lad heat] (15.100,7.340) -- (15.100,4.600);
\draw[lad arr,lad on,draw=ladE] (14.080,3.000) -- node[pos=0.50,rotate=90,lad alab] {ATT} (14.080,1.200);
\draw[lad arr,lad on,draw=ladF] (14.880,4.600) -- node[pos=0.25,rotate=90,lad alab] {FF} (14.880,3.000);
\node[lad axis,anchor=east] at (0.870,1.000) {\ensuremath{-}2};
\node[lad axis,anchor=east] at (0.870,3.000) {0};
\node[lad axis,anchor=east] at (0.870,5.000) {2};
\node[lad axis,anchor=east] at (0.870,7.000) {4};
\node[lad axis,anchor=east] at (0.870,9.000) {6};
\node[lad axis,anchor=east] at (0.870,11.000) {8};
\node[lad axis,anchor=east] at (0.870,13.000) {10};
\node[lad axis,anchor=east] at (0.870,15.000) {12};
\node[lad axis,anchor=east] at (0.870,17.000) {14};
\node[lad axis,anchor=east] at (0.870,18.000) {eV};
\node[lad head,anchor=center] at (4.475,18.000) {Kr and Kr\ensuremath{_2}};
\node[lad head,anchor=center] at (9.650,18.000) {KrF and Kr\ensuremath{_2}F};
\node[lad head,anchor=center] at (14.500,18.000) {F\ensuremath{_2}};
\node[lad lab,anchor=east] at (2.570,3.000) {Kr};
\node[lad lab,anchor=east] at (2.570,12.915) {Kr\ensuremath{^*} 9.915};
\node[lad lab,anchor=east] at (2.570,16.996) {Kr\ensuremath{^+} 13.996};
\node[lad lab2l,anchor=north west] at (5.030,12.150) {Kr\ensuremath{_2^*}\\9.2};
\node[lad lab2,anchor=south east] at (6.300,15.896) {Kr\ensuremath{_2^+}\\12.846};
\node[lad tlab,anchor=west] at (4.460,14.300) {11.30};
\node[lad tlab,anchor=west] at (4.460,14.998) {11.998};
\node[lad tlab,anchor=west] at (4.460,15.750) {12.75};
\node[lad alab,anchor=east] at (2.310,12.195) {G3X};
\node[lad alab,anchor=east] at (5.040,12.557) {AK8\_Kr};
\node[lad hlab,anchor=west] at (5.980,3.370) {0.77};
\node[lad lab,anchor=west] at (12.380,14.396) {Kr\ensuremath{^+} + F\ensuremath{^-} 11.396};
\node[lad lab,anchor=west] at (12.380,13.246) {Kr\ensuremath{_2^+} + F\ensuremath{^-} 10.246};
\node[lad lab,anchor=south east] at (12.300,3.840) {Kr + F 0.8};
\node[lad lab2l,anchor=west] at (11.130,8.856) {KrF(B) 5.856\\KrF(C) 120\,cm\ensuremath{^{-1}} below\\BC, CB mix them};
\node[lad lab2,anchor=east] at (7.220,8.276) {Kr\ensuremath{_2}F\ensuremath{^*}\\5.276};
\node[lad slab,anchor=east] at (8.470,8.566) {G5,\\G5C};
\node[lad hlab,anchor=north] at (9.350,3.420) {X, 0.06};
\node[lad hlab,anchor=north] at (9.750,3.720) {A, 0.50};
\node[lad hlab,anchor=north] at (10.150,3.420) {A, 0.53};
\node[lad alab,anchor=north west] at (10.620,7.991) {QF2, QKR, QE\\QF2C, QKRC, QEC};
\node[lad hlab,anchor=north] at (7.980,3.720) {1.53};
\node[lad alab,anchor=north] at (7.520,7.426) {Q2F};
\node[lad alab,anchor=north] at (8.400,7.426) {Q2FF};
\node[lad lab,anchor=east] at (13.820,3.000) {F\ensuremath{_2}};
\node[lad lab,anchor=east] at (13.820,3.800) {F 0.8};
\node[lad lab,anchor=east] at (13.820,4.600) {F + F 1.6};
\node[lad lab,anchor=east] at (13.820,1.200) {F\ensuremath{^-} + F \ensuremath{-}1.8};
\node[lad lab,anchor=east] at (13.820,0.400) {F\ensuremath{^-} \ensuremath{-}2.6};
\node[lad tlab,anchor=west] at (15.150,6.160) {{\unsansmath\ensuremath{a\,^3\Pi_u}} 3.16};
\node[lad tlab,anchor=west] at (15.150,7.340) {{\unsansmath\ensuremath{A\,^1\Pi_u}} 4.34};
\draw[lad arr,lad on,draw=ladInk2] (1.250,-0.550) -- (1.950,-0.550);
\node[lad lab,anchor=west] at (2.050,-0.550) {a row of the default network};
\draw[lad arr,lad opt,draw=ladInk2] (8.850,-0.550) -- (9.550,-0.550);
\node[lad lab,anchor=west] at (9.650,-0.550) {a row of an optional group, off by default};
\draw[lad arr,lad zero,draw=ladInk2] (1.250,-0.970) -- (1.950,-0.970);
\node[lad lab,anchor=west] at (2.050,-0.970) {a band whose default is 0, no measurement};
\draw[lad hv,lad on,draw=ladB] (8.850,-0.970) -- (9.550,-0.970);
\node[lad lab,anchor=west] at (9.650,-0.970) {a photon, with its row and band};
\draw[lad arr,lad on,draw=ladH] (1.250,-1.390) -- (1.950,-1.390);
\node[lad lab,anchor=west] at (2.050,-1.390) {a quench, the level's energy to heat};
\draw[lad heat] (8.850,-1.390) -- (9.550,-1.390);
\node[lad lab,anchor=west] at (9.650,-1.390) {heat as a lower state falls apart (eV) or cascades};
\draw[lad pair,draw=ladInk2] (1.250,-1.810) -- (1.950,-1.810);
\node[lad lab,anchor=west] at (2.050,-1.810) {a pair, the sum of two formation energies};
\draw[lad tick,draw=ladMuted] (8.850,-1.810) -- (9.550,-1.810);
\node[lad lab,anchor=west] at (9.650,-1.810) {a threshold, or where a photon leaves its lower state};
\end{tikzpicture}}

\makeatletter
\newcommand\wm[1]{\ifcsname wm@#1\endcsname\csname wm@#1\endcsname
  \else\errmessage{no map number '#1'. Rerun make_system_map_v0_2.py}\fi}
\newcommand\wi[1]{\ifcsname wi@#1\endcsname\csname wi@#1\endcsname
  \else\errmessage{no overview number '#1'. Rerun make_overview_v0_2.py}\fi}
\newcommand\lad[1]{\ifcsname lad@#1\endcsname\csname lad@#1\endcsname
  \else\errmessage{no ladder number '#1'. Rerun make_energy_ladder.py}\fi}
\makeatother

% chemistry and units
\newcommand{\unit}[1]{\ensuremath{\,\mathrm{#1}}}
\newcommand{\el}{\ensuremath{\mathrm{e}^-}}
\newcommand{\Krs}{Kr\ensuremath{^*}}
\newcommand{\Krp}{Kr\ensuremath{^+}}
\newcommand{\Krtp}{Kr\ensuremath{_2^+}}
\newcommand{\Krts}{Kr\ensuremath{_2^*}}
\newcommand{\KtF}{Kr\ensuremath{_2}F\ensuremath{^*}}
\newcommand{\Ft}{F\ensuremath{_2}}
\newcommand{\Fm}{F\ensuremath{^-}}
\newcommand{\ra}{\ensuremath{\to}\ }
\newcommand{\hnu}{\ensuremath{h\nu}}
\newcommand{\Edep}{\ensuremath{E_\mathrm{dep}}}
\newcommand{\gIP}{\ensuremath{g_0I_\mathrm{sat}/P_\mathrm{in}}}
\newcommand{\x}[1]{\ensuremath{\times}#1}
% a factor of the measurement as a percentage of it (0.89 -> 89), to the
% factor's own precision: \pctof{x} for two decimals, \pctof[1]{x} for three
\newcommand{\pctof}[2][0]{\pgfmathparse{100*#2}%
  \pgfmathprintnumber[fixed,precision=#1]{\pgfmathresult}}

\begin{document}
\sansmath

{\LARGE\bfseries How accurate LIGHT v0\_2 is}\par
\vspace{4pt}
{\large The basic network of LIGHT v0\_2, version \wm{version}, on 30
September 2026, with its energy ladder, its system map and its reaction
rates}\par
\vspace{10pt}

\noindent\fbox{\parbox{\dimexpr\linewidth-2\fboxsep-2\fboxrule\relax}{%
\setlength{\parskip}{5pt}%
\textbf{Bottom line.} LIGHT v0\_2 reproduces the fitted KrF decays [1, 2] within \wi{sc.train.lo}
to \wi{sc.train.hi}\,\% rms and predicts decays that no fit used [4] within \wi{sc.xu_figs5_7}\,\% rms,
and at the Tahoma laser [5] its default network gives
\pctof{\wi{tah.g0.x}}\,\% of the measured small-signal gain of
\wi{tah.g0.meas}\unit{\%/cm}.

It also gives \pctof{\wi{tah.isat.x}}\,\% of Tahoma's saturation intensity
and \pctof{\wi{tah.eta.x}}\,\% of its formation efficiency. The C/B emission ratio of Yu and Setser [3] is matched to \wi{sc.ys_ir300}\,\%. No fit used Tahoma's numbers and no rate was set by a laser
measurement, so the Tahoma scores are predictions. Most of the gain it misses lies inside the band of
\Krts{} + F \ra KrF(B) + Kr, a rate no one has measured, whose top gives
\pctof{\wi{tah.hiF.x}}\,\%.

Its largest misses are on data that conflict with the fitted sets,
\wi{sc.hakuta_fig2} and \wi{sc.hakuta_fig3}\,\% rms on the C-A decays of
Hakuta et al.\ [6] and \wi{sc.bs_tau_b}\,\% on the KrF(B) lifetime of
Burnham and Searles [7], and the VUV yield of pure krypton at low pump,
\pctof{\wi{pure.vuv.x}}\,\% of Morozov's measurement [8].

{\footnotesize [1] Eden, Waynant, Searles and Burnham, Appl.\ Phys.\
Lett.\ 32, 733 (1978). [2] Gadomski, Xu, Setser and Nelson, Chem.\ Phys.\
Lett.\ 189, 153 (1992). [3] Yu and Setser, J.\ Phys.\ Chem.\ 94, 2934
(1990). [4] Xu, Gadomski and Setser, J.\ Chem.\ Phys.\ 99, 2591 (1993).
[5] Spectra Technology, \emph{Experimental Support for KrF Laser
Modeling}, final report to Los Alamos, 19 July 1989, OSTI 10136715.
[6] Hakuta et al., Jpn.\ J.\ Appl.\ Phys.\ 29, 1194 (1990). [7] Burnham and Searles, J.\ Chem.\ Phys.\ 67, 5967 (1977).
[8] Morozov, Heindl, Kr\"ucken, Ulrich and Wieser, J.\ Appl.\ Phys.\ 103,
103301 (2008).\par}}}
\medskip

LIGHT v0\_2 is v0\_1's working copy. Its input is v0\_1's, a power
deposited uniformly in the electrons, and its basic network adds the
\Krts{} route to v0\_1's, \wi{rates.nBasic} reaction rows in all. Every
rate known only as a published spread, or not at all, is carried as a
band. A band's default is a measurement, or 0 where no measurement exists,
so the default network rests on measurements and theory alone. The
KrF(B,C) block and every formation link are v0\_1's. Every KrF rate in the
block but one is a measurement of that process or a joint fit to kinetic
measurements of it, and the exception is the KrF(B) radiative lifetime,
Kannari's \wm{tauB}\unit{ns} where Eden measured 6.8 $\pm$ 0.2\unit{ns}. No
rate was set by a laser measurement. Scored against Tahoma, the one
independent Kr/F$_2$ laser, under the same beam drive as v0 and v0\_1,
v0\_2's default network gives the small-signal gain at \x{\wi{tah.g0.x}}
the measured value, the saturation intensity at \x{\wi{tah.isat.x}} and
the formation efficiency \gIP{} at \x{\wi{tah.eta.x}}. Kannari's estimates
for the two unmeasured \Krts{} rows give the gain at \x{\wi{tah.kannari.x}},
and \Krts{} + F at the top of its band gives \x{\wi{tah.hiF.x}}. The
largest open questions are the rate of \Krts{} + F \ra KrF(B) + Kr, which
has no measurement, and the trimer rate above \wi{tri.klimek.hi}\unit{K}.
Table~\ref{tab:score} lists every test. The comparison with v0 and v0\_1
is in \texttt{docs/REVIEW/}.

\begin{table}[h]
\centering\footnotesize
\renewcommand{\arraystretch}{1.12}
\begin{tabularx}{\linewidth}{@{}>{\raggedright\arraybackslash}X
  >{\raggedright\arraybackslash}p{2.6cm}>{\raggedright\arraybackslash}p{2.5cm}
  >{\raggedright\arraybackslash}p{2.5cm}@{}}
\toprule
test & measured & v0\_2 & role of the data\\
\midrule
energy ledger closure, \wi{closure.runs} runs across the modes and bands &
  & below \wi{closure.max} of \Edep & bookkeeping\\
KrF(B) decay against Kr and F$_2$, Eden 1978 & & rms \wi{sc.eden_kr} and
  \wi{sc.eden_f2}\,\% & train\\
KrF(B,C) slow decay, Gadomski 1992 Figs. 1 and 2 & & rms \wi{sc.ksu_fig1}
  and \wi{sc.ksu_fig2}\,\% & train\\
C/B emission ratio at 300\unit{K}, Yu and Setser 1990 & \wi{ys.meas}
  $\pm$ \wi{ys.err} & \wi{sc.ys_ir300}\,\% & train\\
KrF(B,C) decays, Xu 1993 Figs. 5 to 7 & & rms \wi{sc.xu_figs5_7}\,\% &
  held out\\
decay constants Xu 1993 printed, 50 and 100 Torr Kr & \wi{xu.50.meas}
  and \wi{xu.100.meas}\unit{s^{-1}} & \x{\wi{xu.50}} and \x{\wi{xu.100}} &
  structure check\\
C-A decays, Hakuta 1990 & & rms \wi{sc.hakuta_fig2} and
  \wi{sc.hakuta_fig3}\,\% & held out, in conflict\\
$\tau_B$, Burnham and Searles 1977 & \wi{sc.bs.meas} $\pm$
  \wi{sc.bs.err}\unit{ns} & \wi{sc.bs_tau_b}\,\% & held out, in conflict\\
C/B ratio 230\unit{K} over 300\unit{K}, Yu and Setser & \wi{ysr.meas}
  $\pm$ \wi{ysr.err} & \wi{sc.ys_ratio_230_300}\,\% & held out, partly
  behind the gap v0\_2 books\\
\midrule
\Krs{} + \Krp{} made per eV, pure Kr and Ar, Saito 2003 & & \x{\wi{pure.Kr.x}}
  and \x{\wi{pure.Ar.x}} & check\\
\Krtp{} + \el{} recombination, Shiu and Biondi 1977 & & \x{\wi{pure.dr.lo}}
  & v0\_2's input\\
\Krp{} + 2Kr \ra \Krtp{} + Kr, Smith et al. 1972 & \wi{pure.conv.lo} to
  \wi{pure.conv.hi} & \wi{pure.k15} & in v0\_2's input\\
VUV yield of pure Kr at low pump, Morozov 2008 & \wi{pure.vuv.meas}
  $\pm$ \wi{pure.vuv.err} & \wi{pure.vuv}, \x{\wi{pure.vuv.x}} & check\\
\quad the same, each band at either end & & moves by \wi{pure.band.move}
  or less & check\\
electron closure at Tahoma, kinetic EEDF loop & $T_e$ \wi{el.loop.te}\unit{eV}
  & \wi{el.light.te}\unit{eV} & reference\\
\midrule
Tahoma small-signal gain $g_0$, v0\_2's default network &
  \wi{tah.g0.meas}\unit{\%/cm} & \wi{tah.g0}, \x{\wi{tah.g0.x}} & device,
  score only\\
\quad the same, Kannari's two \Krts{} estimates & & \wi{tah.kannari.g0},
  \x{\wi{tah.kannari.x}} & device, score only\\
\quad the same, \Krts{} + F at the top of its band & & \wi{tah.hiF.g0},
  \x{\wi{tah.hiF.x}} & device, score only\\
\quad the same, trimer on $T_g^{-1/2}$ & & \wi{tah.tri.g0},
  \x{\wi{tah.tri.x}} & device, score only\\
Tahoma saturation intensity $I_\mathrm{sat}$ & \wi{tah.isat.meas} $\pm$
  \wi{tah.isat.err}\unit{MW/cm^2} & \wi{tah.isat}, \x{\wi{tah.isat.x}} &
  device, score only\\
Tahoma formation efficiency \gIP & \wi{tah.eta.meas} & \wi{tah.eta},
  \x{\wi{tah.eta.x}} & device, score only\\
\bottomrule
\end{tabularx}
\caption{LIGHT v0\_2 against the data, every band at its default unless a
row says otherwise. Rates are in cm$^3$s$^{-1}$ or cm$^6$s$^{-1}$. A
percentage is v0\_2's deviation from the data, signed for a single value
and rms for a set. A \x{} factor is v0\_2 over the measurement.}
\label{tab:score}
\end{table}

\section{How the numbers were made}

Each dataset has one role. Kinetic measurements of the channels being
fitted train them. Other kinetic measurements are held out and test the
fit. Laser performance only scores, and the fitting code refuses a dataset
of that kind. This is the repository's predictive-only policy. Measurements
score the model after the fact and never steer a number, and a band is
never narrowed or moved by a score. Every number in this note is written
by \texttt{docs/make\_overview\_v0\_2.py} and
\texttt{docs/make\_system\_map\_v0\_2.py} from v0\_2 and from \texttt{ML/},
and a test fails when those files no longer match the model.

ML runs v0\_2 when the environment variable \texttt{LIGHT\_VERSION} names
it. v0\_2 takes only a power input. Tahoma was pumped by an electron beam,
so v0\_2 is scored there under the same prototype beam drive as v0 and
v0\_1, a test fixture in \texttt{docs/} and not part of the model. Since
v0\_2's default network carries the \Krts{} route, ML's Tahoma
configuration is that default network, with every band at its default
and $n_e$ free from a seed of \wi{tah.seed}\unit{cm^{-3}}.
Section~\ref{sec:input} runs v0\_2's own input at Tahoma's conditions.
Section~\ref{sec:bands} moves one band at a time to each end of its range,
the others at their defaults.

\section{The KrF(B,C) block}

v0\_2 keeps v0's KrF(B,C) block, as v0\_1 does. KrF has two emitting
states. B lases at 248\unit{nm}, and C radiates slowly in a band at
275\unit{nm}. LIGHT treats them as the two coupled levels of eq. 5 of
Gadomski et al. (1992) and reproduces that model's analytic decay rate to
within \wi{eq5.dev} at every pressure tested. The block has \wi{fit.n}
constants. Nine come from one joint fit to three measurement sets that
agree with one another, the KrF(B) decays of Eden et al. (1978), the Kansas
State B/C decays (Gadomski et al. 1992) and Yu and Setser's equilibrium
C/B ratio (1990). The tenth, $\tau_B$, stays at Kannari's
\wm{tauB}\unit{ns}, inside the fit's interval of \wi{fit.tau_B.lo} to
\wi{fit.tau_B.hi}\unit{ns}. v0\_2 shares it between the 248-nm B-X band and
B-A emission at 272\unit{nm}, which by Dunning and Hay's calculation
carries \wm{pc.BA}\,\% of the B-X rate, so the total radiative rate of B
is unchanged. The fit's median and 68\,\% interval for B-C mixing by Kr is
$k_{BC}=$ \wi{fit.k_BC.mid} [\wi{fit.k_BC.lo}, \wi{fit.k_BC.hi}]
cm$^3$s$^{-1}$. The equilibrium ratio is $K_\mathrm{eq}=$ \wi{fit.K_eq.mid}
[\wi{fit.K_eq.lo}, \wi{fit.K_eq.hi}], and the C lifetime is $\tau_C=$
\wi{fit.tau_C.mid} [\wi{fit.tau_C.lo}, \wi{fit.tau_C.hi}]\unit{ns}. LIGHT
carries the best fit, $k_{BC}=\wm{k.BC}$\unit{cm^3s^{-1}}, $K_\mathrm{eq}=$
\wm{Keq} and $\tau_C=$ \wm{tauC}\unit{ns}. At 1\unit{atm} and 300\unit{K}
the lasing B state holds \wi{fit.bshare.mid} [\wi{fit.bshare.lo},
\wi{fit.bshare.hi}] of the manifold, and LIGHT's rounded constants give
\wi{fit.bshare.light}.

The fitted decay sets are matched to between \wi{sc.train.lo} and
\wi{sc.train.hi}\,\% rms. The held-out decays of Xu et al. (1993) are
matched to \wi{sc.xu_figs5_7}\,\%, and the two decay constants they print
come out at \x{\wi{xu.50}} and \x{\wi{xu.100}}. Hakuta et al. (1990)
measured faster C-A decays, missed by \wi{sc.hakuta_fig2} and
\wi{sc.hakuta_fig3}\,\% rms, and their data conflict with the three fitted
sets. Burnham and Searles (1977) measured $\tau_B=$
\wi{sc.bs.meas}\unit{ns}. Those two conflicts are v0's and remain.

The energy ledger books Yu and Setser's measured B-C gap,
\wm{gapBC}\unit{cm^{-1}}, as v0\_1 does. It sets how $K_\mathrm{eq}$
changes with temperature, and their C/B ratio between 230 and 300\unit{K}
comes out \wi{sc.ys_ratio_230_300}\,\%. That is not an independent test.
Their 120 $\pm$ 50\unit{cm^{-1}} is a consensus of two routes, 110 $\pm$
20 from the absolute ratio and 180 $\pm$ 30 from this very pair of
temperatures.

Tahoma's gain is the sharpest test of the block, since $g_0$ is the
KrF(B) density times the stimulated cross section. Over \wi{post.n} draws
from the fit's posterior the gain is \wi{post.g0.mid}\unit{\%/cm}
[\wi{post.g0.lo}, \wi{post.g0.hi}] against the measured
\wi{tah.g0.meas}.

\section{Charge, recombination and fluorine}

The charge chemistry is v0\_1's, and no electrons are injected. $n_e$
starts from a seed of $\wm{ne}$\unit{cm^{-3}} and follows ionization,
attachment and recombination. At the reference point it runs at
$\wm{neF2}$\unit{cm^{-3}} while F$_2$ lasts and peaks at
$\wm{nePeak}$\unit{cm^{-3}} after F$_2$ burns out at \wm{F2burn}\unit{ns}.
In three-body collisions with Kr, \Krp{} turns into \Krtp{} at
$\wm{k.K15}$\unit{cm^6s^{-1}}, the mean of three measurements. \Krtp{}
recombines with electrons at the rate Shiu and
Biondi measured, $\wm{k.R11p}\,(300\,\mathrm{K}/T_e)^{\wm{pow.R11p}}$
cm$^3$s$^{-1}$. Each recombination returns a \Krs{} and heats the gas. Of
the ions that recombine, \wm{ion.R11p}\,\% meet an electron and
\wm{ion.ii}\,\% meet F$^-$ and form KrF(B), which makes \wm{frm.ii}\,\% of
the KrF formed. Once the F$_2$ is gone, Penning ionization of \Krs{} and
recombination hand the same energy back and forth, so the flow through
\Krs{} passes 100\,\% of \Edep, \wm{krs.in}\,\% here. In v0\_2 part of that
\Krs{} forms \Krts{} and radiates at 147\unit{nm} instead
(Section~\ref{sec:kr2}), and the gas ends at \wm{TgEnd}\unit{K}.

F + F + Kr gives F atoms back to F$_2$ at $\wm{k.FF}$\unit{cm^6s^{-1}},
Cohen and Bott's evaluated Ar value, since no Kr measurement exists. At
the reference point it returns \wm{FF.in}\,\% of \Edep{} from the F atoms
to heat, far too slowly to refill the F$_2$.

\section{The \Krts{} route}
\label{sec:kr2}

v0\_2's basic network carries the \Krts{} excimer, which v0\_1 had as an
optional group. \Krs{} + 2Kr \ra \Krts{} + Kr runs at Ueda's
$\wm{k.AK8}$\unit{cm^6s^{-1}}, measured at 0.5 to 1\unit{atm}, on Wieme and
Lenaerts' $(T_g/300)^{\wm{tpow.AK8}}$, measured over 300 to 451\unit{K}.
\Krts{} radiates at 147\unit{nm} with the measured \wm{tauRAD}-ns
lifetime. \Krts{} + \Ft{} runs at Quigley and Hughes'
$\wm{k.kr2F2}$\unit{cm^3s^{-1}}, split evenly between KrF(B) and \KtF{} by
Kannari's assignment, since the products were not measured. \Krts{} + F
\ra KrF(B) + Kr and \Krts{} + \Krts{} \ra \Krtp{} + 2Kr + \el{} have no
measurement and run at 0. At the reference point \wm{brs.AK8}\,\% of the
\Krs{} forms \Krts, \wm{brs.AK8.F2}\,\% of the \Krs{} used while the F$_2$
lasts. \wm{brk2.RAD}\,\% of that \Krts{} radiates, so \wm{out.147}\,\% of
\Edep{} leaves at 147\unit{nm}, and \Krts{} + \Ft{} forms \wm{frm.kr2}\,\%
of the KrF.

\Krts{} + F decides how much of the route reaches the laser line. With no
measurement its band runs from 0 to \wi{kr2.hi.k}\unit{cm^3s^{-1}}, the
top of a conditional range that SRI gave for the argon analog. At the
reference point the 248-nm share rises from \wi{kr2.def}\,\% at 0 to
\wi{kr2.lo}\,\% at \wi{kr2.lo.k}, \wi{kr2.kann}\,\% at Kannari's
\wi{kr2.kann.k} and \wi{kr2.hi}\,\% at the top, where the 147-nm share
falls from \wi{kr2.vuv} to \wi{kr2.hi.vuv}\,\%. The share saturates
because the F$_2$ burns out at \wi{kr2.burn}\unit{ns} and leaves
\wi{kr2.F}\unit{cm^{-3}} F atoms, which turn \Krts{} into KrF
\wi{kr2.kann.ratio} times faster than \Krts{} radiates at Kannari's rate
and \wi{kr2.hi.ratio} times faster at the top, and each KrF photon gives
its F atom back. At the top \wi{kr2.hi.afterpc}\,\% of the gain comes
after the burnout. With Kannari's estimates for both unmeasured rows,
\wi{kr2.kann.k} and \wi{kr2.pen.k}\unit{cm^3s^{-1}}, v0\_2 gives v0\_1's
optional \Krts{} group, \wi{kr2.both}\,\%.

Pure krypton tests the route without F$_2$. At Morozov's low pump v0\_2
gives a VUV yield of \wi{pure.vuv}, where Morozov et al. (2008) measured
\wi{pure.vuv.meas} $\pm$ \wi{pure.vuv.err}, and no band at either end
moves it by more than \wi{pure.band.move}. So the miss does not come from
the route, and that point cannot score the bands. At
\wi{pure.hi.P}\unit{kW\,cm^{-3}} for \wi{pure.hi.pulse}\unit{\mu s} in
\wi{pure.hi.torr}\unit{Torr} of Kr the bands move the yield from
\wi{pure.hi.lo}, with \wi{pure.hi.lo.what}, through \wi{pure.hi.def} at
the defaults to \wi{pure.hi.hi}, with \wi{pure.hi.hi.what}. A pure-Kr VUV
yield measured against pump rate would score them, and none is among the
sources held.

\section{The bands}
\label{sec:bands}

v0\_2 carries \wm{bd.n} rates as bands. A measured rate keeps its
measurement as the default and spans the other published values. A rate
with no measurement is 0 by default and runs to the largest published
value, which is no bound, and \wm{bd.nzero} bands are of that kind.
Analogs and device data never set a default or narrow a band.
Appendix~A.4 gives each band's sources. Table~\ref{tab:bands} moves one
band at a time to each end at the reference point. Two bands move the
248-nm share far, \Krts{} + F from \wm{bd.kr2_f.lo248} to
\wm{bd.kr2_f.hi248}\,\% and the trimer's temperature law from
\wm{bd.trimer_tg_pow.lo248} to \wm{bd.trimer_tg_pow.hi248}\,\%. Each of
the other six moves it by \wm{bd.minor} points or less.

\begin{table}[H]
\centering\footnotesize
\renewcommand{\arraystretch}{1.12}
\setlength{\tabcolsep}{4pt}
\begin{tabularx}{\linewidth}{@{}>{\raggedright\arraybackslash}X
  >{\raggedright\arraybackslash}p{1.45cm}>{\raggedright\arraybackslash}p{3.6cm}
  >{\raggedright\arraybackslash}p{1.65cm}>{\raggedright\arraybackslash}p{1.7cm}
  >{\raggedright\arraybackslash}p{1.55cm}@{}}
\toprule
band & default & low to high end & 248\unit{nm} & 147\unit{nm} &
  $T_g$ at the end (K)\\
\midrule
\wmBands
\bottomrule
\end{tabularx}
\caption{Each band alone at its two ends at the reference point, the
others at their defaults. The shares are percent of \Edep. With every band
at its default v0\_2 sends \wm{out.248}\,\% out at 248\unit{nm} and
\wm{out.147}\,\% at 147\unit{nm}, and the gas ends at \wm{TgEnd}\unit{K}.}
\label{tab:bands}
\end{table}

Tahoma scores the bands and sets none (Table~\ref{tab:tahbands}). Its gain
moves most with the trimer's weak law, to \x{\wi{tah.tri.x}}, so Tahoma
counts against that law. \Krts{} + F at the top of its band brings the
gain to \x{\wi{tah.hiF.x}}. Each of the other six moves it by
\wi{tah.bd.minor}\unit{\%/cm} or less. Tahoma's window, \wi{tah.win.lo}
to \wi{tah.win.hi}\unit{ns}, closes with \wi{tah.f2left}\,\% of the F$_2$
fill left, while at the reference point \wi{kr2.hi.afterpc}\,\% of the
band's gain comes after the F$_2$ has burnt out. So \Krts{} + F moves
Tahoma's gain far less than the reference point's 248-nm share.

\begin{table}[H]
\centering\footnotesize
\renewcommand{\arraystretch}{1.12}
\begin{tabularx}{0.82\linewidth}{@{}>{\raggedright\arraybackslash}X
  >{\raggedright\arraybackslash}p{2.6cm}>{\raggedright\arraybackslash}p{2.9cm}@{}}
\toprule
band & $g_0$ (\%/cm), low to high end & over the measured
  \wi{tah.g0.meas}\\
\midrule
\wiTahBands
\bottomrule
\end{tabularx}
\caption{Each band alone at its two ends at Tahoma, under the beam drive
with ML's constants, the others at their defaults. With every band at its
default the gain is \wi{tah.def.g0}\unit{\%/cm}.}
\label{tab:tahbands}
\end{table}

\section{KrF formation}

The chain from \Krp{} and \Krs{} to KrF is v0\_1's. \Krs{} + \Ft{} reacts
at $\wm{k.harp}$\unit{cm^3s^{-1}}, the mean of three first-hand
measurements. Velazco, Kolts and Setser (1976) measured the KrF$^*$ yield
of that reaction as \wm{pc.harpY}\,\%, and Yu and Setser (1990) put
\wm{pc.fC}\,\% of that KrF$^*$ in C. Tahoma measured \gIP{} $=$
\wi{tah.eta.meas}. ML puts a 248-nm probe flux on v0\_2 and fits its
$I_\mathrm{sat}$ at fixed $g_0$, as Tahoma did. That gives \wi{tah.eta}.
Its KrF formation rate in 5-eV quanta is \wi{tah.form} of the pump. The
difference is KrF that forms in C, or mixes into it, and decays there
without meeting the laser field.

Table~\ref{tab:links} moves one link at a time to another value it could
take. The measured range of the harpoon yield, 0.70 to 1.0, moves \gIP{}
from \wi{link.yield.lo} to \wi{link.yield.hi} and brackets Tahoma's
\wi{tah.eta.meas}. Three links have no measurement at all. They are the
KrF$^*$ yield of ion-ion recombination with F$^-$, the share of that KrF
formed in C, and \Krts{} + F \ra KrF. Only analogs exist for the first
two, and both analogs lower \gIP{} toward the measurement. v0\_2 runs the
third at 0, and Kannari's estimate or the top of its band raise \gIP{}.
The predictive-only policy lets no analog set a link and no link be chosen
by its fit to Tahoma. The rest of the \Krts{} route moves \gIP{} by
\wi{link.minor} at most.

\begin{table}[H]
\centering\footnotesize
\renewcommand{\arraystretch}{1.12}
\begin{tabularx}{\linewidth}{@{}>{\raggedright\arraybackslash}X
  >{\raggedright\arraybackslash}p{1.8cm}>{\raggedright\arraybackslash}p{1.9cm}
  >{\raggedright\arraybackslash}p{3.3cm}r r@{}}
\toprule
link & v0\_2 & tried & what the value is & \gIP & change\\
\midrule
\wiLinks
\bottomrule
\end{tabularx}
\caption{Each formation link moved alone from v0\_2's chain at Tahoma,
\wi{tah.f2torr}\unit{Torr} of F$_2$ in Kr at \wi{tah.torr}\unit{Torr} in
all, pumped at \wi{tah.pin}\unit{kW\,cm^{-3}}. v0\_2's chain gives
\wi{link.base}, and Tahoma measured \wi{tah.eta.meas}.}
\label{tab:links}
\end{table}

\section{The electron side}

v0\_2 integrates each krypton electron-impact rate over v0's EEDF with
two temperatures, a Maxwellian bulk at $T_e$ and a tail above
9.915\unit{eV} read from kinetic EEDFs. Each captured electron takes the
energy it really carries. For \Ft{} attachment at Tahoma that is
\wi{el.eatt}\unit{eV}, where $\tfrac32 T_e$ would book
\wi{el.eatt.32}\unit{eV}. \Krtp{} recombination goes as
$T_e^{-\wi{el.rpow}}$, so a recombining electron takes
\wi{el.rcap}\,$T_e$. A kinetic loop that iterates v0\_2 with a Boltzmann
solution, with the same recombination law as its sink, is the reference.
At Tahoma v0\_2's own closure lands within \wi{el.dev.te}\,\% of the
loop's $T_e$ and \wi{el.dev.ne}\,\% of its $n_e$, and within
\wi{el.dev.g0} and \wi{el.dev.formation}\,\% of its gain and formation.

The cross sections are v0's. With Biagi's krypton set the computed $W$ is
\wi{sw.krw.lo} to \wi{sw.krw.hi} times Nakamura's measured values. For
\Ft{} attachment the groups that measured the rate constant disagree
beyond their errors, and SIGLO gives between \wi{sw.f2.lo} and
\wi{sw.f2.hi} times their values.

On pure gases the beam drive's Spencer-Fano yields make \wi{pure.Kr.x}
times the \Krs{} and \Krp{} per eV that the measured $W$ and
$N_\mathrm{ex}/N_i$ give in Kr, and \wi{pure.Ar.x} times in Ar. v0\_2
takes its \Krtp{} recombination from Shiu and Biondi and its \Krp{}
conversion from a mean that includes Smith et al., so two of v0's
pure-gas checks are v0\_2's inputs. The VUV yield of pure Kr at low pump
is still the one miss. Without the thermal electrons' excitation v0\_2
gives \wi{pure.vuv.nothermal}, close to the one photon per excitation and
per ion that the measured $W$ and $N_\mathrm{ex}/N_i$ allow,
\wi{pure.vuv.ideal}, so the shortfall is either physics LIGHT lacks or a
systematic error of the measurement.

\section{The power input}
\label{sec:input}

v0\_2 feeds its power to the electron gas. Run that way at Tahoma's
conditions, with no beam, it gives $g_0=$ \wi{tah.own.g0}\unit{\%/cm},
\x{\wi{tah.own.x}} the measured value. Its electrons then carry the power
at $T_e=$ \wi{tah.own.te}\unit{eV} and $n_e=$ \wi{tah.own.ne}\unit{cm^{-3}},
where the beam-sustained plasma sits at \wi{tah.beam.te}\unit{eV} and
\wi{tah.beam.ne}\unit{cm^{-3}}. The gain comes near, but the electron gas
does not resemble the beam's, so the agreement does not test the power
input. No discharge-pumped Kr/F$_2$ measurement has been scored.

\section{The energy ladder}

Figure~\ref{fig:ladder} places every species of v0\_2 at its formation
energy, the energy above ground-state Kr and \Ft{} at rest that
Appendix~A lists, and draws the rows that move energy between them. It
follows the ladder view of the website's excimer simulation, with
v0\_2's own numbers. KrF(B) lies \lad{DeB}\unit{eV} below the ion pair \Krp{} + \Fm{}, the
measured $D_e(B)$, and \lad{BoverKrF}\unit{eV} above Kr + F. The
\lad{wl.G4}-nm photon carries \lad{hv.G4}\unit{eV} of that and leaves
KrF(X) only \lad{heat.G4}\unit{eV} above Kr + F. The harpoon G3 heats the
gas by \lad{heat.G3}\unit{eV} for each KrF(B) it makes, the ion-ion
recombination K26 by \lad{heat.K26}\unit{eV} and the trimer's formation by
\lad{heat.G5}\unit{eV}. \KtF{} emits \lad{hv.C3}\unit{eV} at
\lad{wl.C3}\unit{nm} and leaves \lad{heat.C3}\unit{eV} of lower-state
repulsion. On the \Krts{} route, which is basic in v0\_2, \Krts{} emits
\lad{hv.V}\unit{eV} at \lad{wl.V}\unit{nm} and leaves
\lad{heat.V}\unit{eV}, or meets F$_2$ and makes KrF(B) through NK7 or
\KtF{} through K2F\_BR.

\begin{figure}[p]
\centering
\ladpicture
\caption{The energy ladder of LIGHT v0\_2, version \lad{version}, in eV
above ground-state Kr and \Ft{} at rest. A bar is a species at its
formation energy and a dashed bar a pair, the sum of two. An arrow is a
row, labelled with its id. It is solid in the default network and dotted
for a band whose default is 0, \lad{zero}. The rows of the optional group
\lad{groups} are dashed. A photon's arrow drops by its quantum, and the red
squiggle below it is the heat its lower state gives the gas as it falls
apart. A red stub is a quench. The energy a row releases heats the gas,
except a photon's, which leaves, and that of \lad{toe}, which goes to the
electrons. KrF(B) and KrF(C) lie \lad{gapBC}\unit{cm^{-1}} apart and are
drawn as a pair of bars that every B and C row meets. The rows
\lad{notdrawn} move nothing between levels and are not drawn.}
\label{fig:ladder}
\end{figure}

\section{The system map}

The map on the next page draws v0\_2's basic network at the notes'
reference point, \wm{p}\unit{atm} of Kr with \wm{xF2}\,\% F$_2$,
\wm{P}\unit{kW\,cm^{-3}} for \wm{pulse}\unit{\mu s}, with $n_e$ seeded at
$\wm{ne}$\unit{cm^{-3}} and then free and every band at its default. Band
widths are proportional to the energy they carry, and every box conserves
it (\wm{nodes} balances, each to $10^{-8}$ of \Edep). Of the deposited
energy, \wm{krf.eff}\,\% passes through KrF in 5-eV quanta. Of that KrF,
\wm{brf.G4}\,\% radiates at 248\unit{nm}, \wm{brf.tri}\,\% forms the
trimer and \wm{brf.e}\,\% is quenched by electrons, so \wm{out.248}\,\% of
the deposited energy leaves at 248\unit{nm}, and \wm{out.147}\,\% leaves
at 147\unit{nm} from \Krts.

% the ladder is a float, which pdfpages would leave pending past the map
\clearpage
\includepdf[pages=1,fitpaper=true]{system_map_v0_2.pdf}

\section{Where v0\_2 is not yet reliable}

\begin{itemize}
\item \Krts{} + F \ra KrF(B) + Kr. No measurement exists, so v0\_2 runs it
at 0 and carries it as a band up to \wi{kr2.hi.k}\unit{cm^3s^{-1}}.
Kannari's \wi{kr2.kann.k} traces to an SRI report whose model leaves the
reaction out and whose one number is a conditional range for Ar$_2^*$ + F.
At the reference point the band spans \wm{bd.kr2_f.lo248} to
\wm{bd.kr2_f.hi248}\,\% at 248\unit{nm}, and at Tahoma \wi{tah.bd.kr2_f.lo}
to \wi{tah.bd.kr2_f.hi}\unit{\%/cm}. A measurement or a bound between
$10^{-12}$ and $10^{-10}$\unit{cm^3s^{-1}} would matter most.
\item The trimer above \wi{tri.klimek.hi}\unit{K}. KrF + 2Kr \ra \KtF{}
takes \wm{brf.tri}\,\% of the KrF at the reference point, where the gas
ends at \wm{TgEnd}\unit{K}. The $(T_g/300)^{\wm{tpow.G5}}$ law is Shui's
calculation, measured only up to \wi{tri.klimek.hi}\unit{K}. With the
pure-Kr $T_g^{-\wi{tri.lo.pow}}$ law, the weakest the data allow, the
reference point sends \wm{bd.trimer_tg_pow.hi248}\,\% out at 248\unit{nm}
instead of \wm{out.248}\,\%. Tahoma's window sits at \wi{tah.tg}\unit{K},
and there that law gives the gain at \x{\wi{tah.tri.x}} the measured
value. Tahoma's gain counts against the weak law as a score and sets
nothing.
\item \Krs{} + 2Kr above 451\unit{K}. Its temperature law is measured over
300 to 451\unit{K}, and the reference point ends at \wm{TgEnd}\unit{K}.
The report warns when the gas passes it.
\item Hot gas. Recombination returns the ionization energy as heat. At
\wi{hot.p}\unit{atm} and \wi{hot.P}\unit{kW\,cm^{-3}} the gas reaches
\wi{hot.tg}\unit{K}, beyond every rate law's data, and \KtF{} and \Krtp{}
redissociation are not in the model. The report warns above 1000\unit{K}.
\item Ion-ion recombination with F$^-$. Its rate was measured only below
300\unit{Torr}, where Jacob et al. (1978) fitted KrF$^*$ pulse shapes.
Rokni et al. (1977) found that \Krtp{} + F$^-$ forms KrF$^*$ as
efficiently as \Krp{} + F$^-$, but from relative intensities. The absolute
KrF$^*$ yield, its B/C split and the rate near 1\unit{atm} have no
measurement. Table~\ref{tab:links} shows what the analogs would do.
\item \Krts{} + \Krts{} and \Krs{} + \Krs{}. Neither Penning rate has a
measurement among the sources held, and at high pump they set the pure-Kr
VUV yield (Section~\ref{sec:kr2}).
\item \Krts's well. REBUILD's estimate puts \Krts{}
\wm{E.kr2s.bind}\unit{eV} below \Krs{} + Kr, since no measured or
calculated binding energy is among the sources held. It only moves energy
between heat and storage.
\item F + F + Kr is held at its 300\unit{K} value, which overstates it in
hot gas.
\item Hakuta's C-A decays and Burnham and Searles' $\tau_B$ remain
unexplained conflicts.
\end{itemize}

\appendix
\section{Reaction rates}

Rates are the values at 300\unit{K}. The column $T$ gives the law v0\_2
applies to them. Electron-impact rows use $\langle\sigma v\rangle$ over the
EEDF, and Section A.3 lists their Maxwellian values. A row that reads a
band prints the value v0\_2 runs by default and the band's name, and
Section A.4 lists the bands. Each row's full citation is in
\texttt{v0\_2/network.py}. The row ids are the ones the code and the
system map use.

\footnotesize
\setlength{\tabcolsep}{3.5pt}
\renewcommand{\arraystretch}{1.12}

\subsection{The basic network, \wi{rates.nBasic} rows}
\wi{rates.nBandRows} of them read a band.
\begin{longtable}{@{}l>{\raggedright}p{4.1cm}>{\raggedright}p{3.1cm}
  >{\raggedright}p{1.9cm}>{\raggedright\arraybackslash}p{4.75cm}@{}}
\toprule
row & reaction & rate & $T$ & source\\
\midrule
\endhead
\bottomrule
\endfoot
\wirBasic
\end{longtable}

\subsection{Optional group on the basic network, \wi{rates.nGroups} rows}
The group \wi{rates.groups} is off at the reference point and switches on
with \texttt{-{}-add}.
\begin{longtable}{@{}l>{\raggedright}p{4.1cm}>{\raggedright}p{3.1cm}
  >{\raggedright}p{1.9cm}>{\raggedright\arraybackslash}p{4.75cm}@{}}
\toprule
row & reaction & rate & $T$ & source\\
\midrule
\endhead
\bottomrule
\endfoot
\wirGroups
\end{longtable}

\subsection{Electron-impact rate coefficients}
Maxwellian $\langle\sigma v\rangle$ in cm$^3$s$^{-1}$ at the $T_e$ of each
column. The last column is the factor the two-temperature tail applies at
the reference point's mid-pulse state, $T_e=$ \wi{rates.teMid}\unit{eV}
with the tail at \wi{rates.rMid}\,$T_e$. The rows below the rule give the
energy an attached electron takes from the electron gas.
\begin{center}\scriptsize
\begin{tabular}{@{}llrrrrrr@{}}
\toprule
row & process & \multicolumn{5}{c}{$T_e$ (eV)} & tail\\
\cmidrule(lr){3-7}
 & & \wirTe & factor\\
\midrule
\wirXsec
\bottomrule
\end{tabular}
\end{center}

\subsection{The bands}
\texttt{network.BANDS}. \texttt{Params.rates} sets a band, and
\texttt{Params.trimer\_tg\_pow} the trimer law.
\begin{longtable}{@{}l>{\raggedright}p{3.2cm}>{\raggedright}p{1.6cm}
  >{\raggedright}p{2.6cm}>{\raggedright\arraybackslash}p{5.2cm}@{}}
\toprule
band & process & default & range & source\\
\midrule
\endhead
\bottomrule
\endfoot
\wirBands
\end{longtable}

\subsection{Where v0\_2's v0 network differs}
The v0 network keeps the notes' six reactions with REBUILD's values, and
its optional groups, \wi{rates.v0groups}, are REBUILD's. These rows differ
from v0\_2's basic network or exist only in the v0 network. The rows
ending in \_C are the bc group's copies for KrF(C).
\begin{longtable}{@{}l>{\raggedright}p{4.1cm}>{\raggedright}p{3.1cm}
  >{\raggedright}p{1.9cm}>{\raggedright\arraybackslash}p{4.75cm}@{}}
\toprule
row & reaction & rate & $T$ & source\\
\midrule
\endhead
\bottomrule
\endfoot
\wirVzero
\end{longtable}

\needspace{20\baselineskip}
\subsection{Formation energies}
Energies above ground-state Kr and F$_2$ at rest. Every reaction's energy
release is the difference of these numbers, so no row can make or lose
energy. $c_v$ is the heat capacity per particle in units of $k_B$.
\begin{center}
\begin{tabular}{@{}lrrl@{}}
\toprule
species & $E_\mathrm{form}$ (eV) & $c_v$ & what sets it\\
\midrule
\wirSpecies
\bottomrule
\end{tabular}
\end{center}

\end{document}
